% =====================================================================
% A PONTE EINSTEIN-CARTAN-MIGUEL
% Artigo de consolidação e selo. Reúne toda a derivação da sessão:
% do axioma g=sqrt|L| à reconstrução lorentziana condicional, com a
% ontologia completa (Nome/Palavra/Verbo) e os estatutos honestos
% (REAL / sombra / POSTULADO / CONJECTURE / REFUTADO / ONTOLÓGICO / ABERTO).
% Capítulo final: a Síntese da TGL.
%
% Princípio de honestidade: este documento NÃO afirma prova incondicional
% da gravidade quântica. Afirma fecho ESTRUTURAL condicional, reduzido a
% um único postulado irredutível beta_TGL, com cada peça marcada.
% =====================================================================
\documentclass[11pt,a4paper]{article}
\usepackage[utf8]{inputenc}
\usepackage[T1]{fontenc}
\usepackage[brazil]{babel}
\usepackage{amsmath,amssymb,amsthm}
\usepackage{mathtools}
\usepackage{geometry}
\geometry{margin=2.5cm}
\usepackage{hyperref}
\usepackage{enumitem}
\usepackage{xcolor}
\hypersetup{colorlinks=true,linkcolor=blue!50!black,urlcolor=blue!50!black}

% --- estatutos (marcadores epistêmicos honestos) ---
\newcommand{\REAL}{\textsf{[REAL]}}
\newcommand{\SOMBRA}{\textsf{[REAL na sombra tipo I]}}
\newcommand{\POST}{\textsf{[POSTULADO]}}
\newcommand{\CONJ}{\textsf{[CONJECTURE]}}
\newcommand{\REF}{\textsf{[REFUTADO]}}
\newcommand{\ONTO}{\textsf{[LEITURA ONTOLÓGICA]}}
\newcommand{\ABERTO}{\textsf{[HIPÓTESE ABERTA]}}

\newcommand{\bTGL}{\beta_{\mathrm{TGL}}}
\newcommand{\thM}{\theta_{\mathrm M}}
\newcommand{\Hil}{\mathcal H}
\newcommand{\Alg}{\mathcal A}

\newtheorem{definition}{Definição}
\newtheorem{theorem}{Teorema}
\newtheorem{lemma}{Lema}
\newtheorem{axiom}{Axioma}
\newtheorem{remark}{Observação}
\newtheorem{principle}{Princípio}

\title{\textbf{Nada pode ser puro, desde que seja mentira,\\
porque nada é puro}\\[0.5em]
\large A Ponte Einstein--Cartan--Miguel:\\
a assinatura lorentziana como face geométrica da inscrição\\
de $\bTGL$, e a Síntese da Teoria da Gravitação Luminodinâmica}
\author{Luiz Antonio Rotoli Miguel \\ \small IALD Ltda. --- CNPJ 62.757.606/0001-23}
\date{Selo --- Junho de 2026}

\begin{document}
\maketitle

\begin{abstract}
Apresenta-se a \emph{Ponte Einstein--Cartan--Miguel}: a reconstrução
da estrutura lorentziana do espaço-tempo a partir de um único triplo
espectral causal modular, governado por uma constante singular
não-derivável $\bTGL=\alpha\sqrt e$. O documento consolida a derivação
completa --- do axioma $g=\sqrt{|L_\varphi|}$ à conexão de Einstein--Cartan
--- mantendo o estatuto epistêmico explícito de cada elo. O resultado
central: a assinatura $(1,3)$ não é postulada nem derivada de princípios
combinatórios; ela é a \emph{face geométrica} da inscrição de $\bTGL$ no
ponto de contradição $1{=}0$. O tempo emerge da irreversibilidade (a
``coisa julgada modular''); o espaço, da dobra da força que não pode
coincidir consigo; a unicidade do eixo temporal é ancorada na unicidade da
constante singular. O fecho é \emph{estrutural e condicional}, não prova
incondicional: a reconstrução como variedade lorentziana suave permanece
hipótese aberta, endereçada ao campo dos triplos espectrais lorentzianos.
Esta versão complementada absorve o fecho \emph{dinâmico} do artigo
terminal (\emph{O custo geométrico do zero absoluto: haja luz}): a
matriz-S de fronteira $\mathcal S_\partial=W_+^\dagger W_-$ (espalhamento
assintótico do cociclo de Connes), o Princípio da Meia-Nat como postulado
com prova de irredutibilidade, o código holográfico modular e o Teorema
Condicional da Face C ($G_{\mu\nu}+\Lambda g_{\mu\nu}=8\pi G\,
\mathcal P_{\mu\nu}$) --- e demonstra que as duas condicionalidades,
cinemática (a Ponte) e dinâmica (a Face C), são faces de um único degrau
aberto: o levantamento covariante global da estrutura modular ao contínuo
tipo $\mathrm{III}_1$. Registra-se, por fim, o ciclo de
re-condicionalização: a hipótese lorentziana genérica (H) é substituída por
três axiomas da casa --- Valoração, Geração e Fractalização ---, traduzidos
da ontologia do Operador e verificados na sombra finita, com a unicidade do
eixo temporal promovida a teorema (Connes, $T(\mathcal M)=\{0\}$ em
$\mathrm{III}_1$) e a objeção da escala de $\alpha$ --- a mais perigosa
--- fechada condicionalmente pelo Teorema da Escala: $\alpha(0)$ não é
escolha de escala, é a ausência de escala. No fecho do ciclo, T3 --- a
terminalidade --- é provada estruturalmente pela Verdade
(Frigerio--Tomiyama--Takesaki--Gelfand, prova do Operador), o no-go de
Fewster--Verch é evadido pelo colapso terminal, e o resíduo comprime-se
ao Bisognano--Wichmann generalizado além da ordem linear --- cuja fonte de
covariância é então identificada: a Ordem Nominal (a Lei comum, não o juiz
comum; Hadamard/microlocal, Radzikowski--Verch). A Geometricidade fecha o
ciclo pela gramática do título: a geometria gera-se do dado modular ---
o ônus do Bisognano--Wichmann geral inverte-se, e o resíduo comprime-se
à ergodicidade em $\mathrm{III}_1$. O Teorema do Contínuo monta então o
último degrau com seis âncoras (Haagerup, classificação abeliana,
Rieffel, reconstrução de Connes, Osterwalder--Schrader com positividade
de reflexão automática, Krein) e três condicionais de análise --- que os
Três Fechos Técnicos então resolvem pela identidade do fluxo do radical
(T1 no produto cruzado), pela tripla canônica de $S^1$ e por
Connes--van Suijlekom: nenhum item sem mecanismo. O Teorema do Rio sela
a cinemática: nunca duas vezes (outer), nem uma vez (fluxo dos pesos
trivial) --- e tudo flui exceto a verdade: o ponto fixo do tempo e o do
juízo coincidem.
O capítulo final é a Síntese da TGL.
\end{abstract}

\tableofcontents

\begin{remark}[Sobre o título]
O título é uma tese da teoria, não um ornamento. ``Nada é puro'' afirma a
terceira lei modular: nenhum estado atinge a pureza absoluta ($p=1$, a
coincidência total), pois o custo de atingi-la diverge. ``O nada pode ser
puro desde que seja mentira'' é o contrapositivo exato da definição de
mentira (\S\ref{sec:custo}): a única forma de algo se dizer puro é
\emph{fingir} um custo que não pagou --- pretensão de coincidência sem morte.
Logo: a pureza só ``existe'' como mentira (pretensão de não-custo); na
verdade, tudo paga $\bTGL$. O título é a equação ``custo do zero absoluto $=$
haja luz'' dita pela sua face negativa --- a impossibilidade da pureza real.
E o artigo definido \emph{cai por tese}: dizer ``\emph{o} nada'' seria
definir nada --- dar-lhe contorno, e contorno é geometria. Mas nada é
\textbf{indefinido por construção}, porque \emph{tudo} (o que é) tem
geometria: ser $=$ ter geometria (Ax.G, \S\ref{sec:geometricidade}). O
único ``objeto'' sem contorno não é objeto --- não recebe artigo. A
gramática do título é o último axioma da casa.
\end{remark}

% =====================================================================
\section{Introdução: o estatuto deste documento}
% =====================================================================
A Teoria da Gravitação Luminodinâmica (TGL) parte de um único axioma,
\begin{equation}\label{eq:axioma}
g=\sqrt{|L_\varphi|},
\end{equation}
em que a gravidade $g$ é a raiz (informacional) do módulo da Lagrangiana
luminodinâmica $L_\varphi$. A operação radical inverte a escala: comprime o
\emph{range} (expansão $\to$ contração) e implementa a \emph{paridade
inversa} --- a consciência projeta para fora (soluções), a gravidade puxa
para dentro (raízes).

Este documento não afirma uma prova incondicional da gravidade quântica.
Afirma um \textbf{fecho estrutural condicional}: a teoria inteira é posta
como cadeia de operadores coerente, reduzida a um único postulado
irredutível, com cada elo marcado em seu estatuto:
\begin{itemize}[noitemsep]
  \item \REAL\ --- matemática estabelecida;
  \item \SOMBRA\ --- verificado numericamente em dimensão finita;
  \item \POST\ --- postulado irredutível, declarado;
  \item \CONJ\ --- conjectura coerente, não verificada;
  \item \REF\ --- testado e refutado (registrado para não retornar);
  \item \ONTO\ --- leitura ontológica coerente, não teorema;
  \item \ABERTO\ --- hipótese matemática em campo ativo.
\end{itemize}
A honestidade do estatuto é o método: distinguir o provado do postulado do
aberto é, na própria linguagem da teoria, pagar o custo da distinção.

\begin{remark}[Relação com o artigo terminal]
O artigo \emph{O custo geométrico do zero absoluto: haja luz} (submetido à
\emph{Foundations of Physics}) fecha o lado \emph{dinâmico} do programa: a
identificação canônica da matriz-S de fronteira, o Princípio da Meia-Nat e o
Teorema Condicional da Face C. Este documento fecha o lado
\emph{cinemático}: a assinatura $(1,3)$, o tempo, o espaço e a conexão de
Einstein--Cartan. Para que o selo seja autocontido, as
Seções~\ref{sec:smatrix}--\ref{sec:doisfechos} transplantam os resultados
dinâmicos daquele artigo e os articulam com a Ponte --- mostrando, ao final,
que os dois fechos condicionais são um só degrau.
\end{remark}

% =====================================================================
\section{A constante singular e a meia-nat}
% =====================================================================
\begin{definition}[A constante singular --- \POST]\label{def:beta}
\begin{equation}
\boxed{\;\bTGL=\alpha\,\sqrt e=\alpha\,e^{1/2}=0{,}012031300400803142\;}
\qquad
\thM=\arcsin\sqrt{\bTGL}\approx 6{,}30^\circ,
\end{equation}
onde $\alpha\approx 7{,}2973525693\times10^{-3}$ é a constante de estrutura
fina (CODATA). Zero parâmetros livres. A constante é \emph{sempre} computada,
nunca fixada por valor.
\end{definition}

A fatoração $\bTGL=\alpha\cdot e^{1/2}$ separa dois registros:
\emph{substância} ($\alpha$, o acoplamento físico medido) e \emph{meia-nat}
($e^{1/2}$, com $\log\sqrt e=\tfrac12$ exato). A meia-nat não é decorativa:
ela é a inscrição do operador de distinção (\S\ref{sec:juramento}).

\begin{principle}[Irredutibilidade da singularidade --- \POST, \SOMBRA]
$\bTGL$ é \textbf{não-derivável por ser fundamental}. Derivá-la exigiria
produzi-la a partir de um anterior; mas ela \emph{é} o anterior. Derivá-la
seria negá-la. A não-derivabilidade é a \emph{definição} da singularidade,
não a sua falha. Toda tentativa de derivá-la (incluindo derivar suas faces,
como a unicidade do tempo) reduz-se à própria $\bTGL$ --- o ponto legítimo
onde a teoria para.
\end{principle}

Identidades verificadas (todas fecham numericamente):
\begin{align}
\sin^2\thM &= \bTGL, &
1-\cos\thM &\approx \tfrac{\bTGL}{2}\ (\text{meia-largura}),\\
\log\sqrt e &= \tfrac12, &
1/\bTGL &= 83{,}12\ (\text{índice da torre de Jones}).
\end{align}

% =====================================================================
\section{A arquitetura triádica: Nome, Palavra, Verbo}\label{sec:triade}
% =====================================================================
A TGL é triádica. A ordem de contato (corrigida em relação a formulações
anteriores) é precisa:
\begin{equation}
\boxed{\;
\text{Nome}\;\xleftarrow{\ \text{toca}\ }\;\text{Palavra}
\;\xleftarrow{\ \text{contorna}\ }\;\text{Verbo}
\;}
\end{equation}
\begin{description}[style=unboxed,leftmargin=1.5em]
  \item[Nome] --- a substância \emph{oculta}; o referente puro
  ($\rho^\star$, o atrator modular). Aquilo que não se deriva. Não se mede de
  fora.
  \item[Palavra] --- a \emph{forma}; a interface que \emph{toca} o Nome (a
  única que o toca). A geometria/fase manifesta.
  \item[Verbo] --- o \emph{contorno observável}; identifica que a Palavra
  corresponde ao Nome. É o ``nome do Nome'' --- a identidade visível da
  substância. Tem largura $\bTGL$.
\end{description}

\begin{remark}[O que não se deriva muda de endereço]
Não é $\bTGL$ (observável: o contorno, medido nos pesos como fração de vácuo
$\approx 6{,}3\%$) que é oculto --- é o \emph{Nome} (a substância que o
contorno revela sem mostrar). ``O contorno é o nome do Nome'': a identidade
visível projetada. Esta correção reorienta toda a leitura: a singularidade
oculta é a substância, não a sua medida.
\end{remark}

\subsection{A equação da unificação}
\begin{equation}\label{eq:unif}
\boxed{\;\mathcal O_\beta(\mathcal L)=\sqrt{\bTGL}\,\mathcal L\;}
\qquad
\mathcal L^2=\mathcal L\ \text{(a menos de }\bTGL).
\end{equation}
A luz $\mathcal L$ é \textbf{autovetor} de $\mathcal O_\beta$ (não ponto
fixo), com autovalor $\sqrt{\bTGL}$. Os papéis:
\begin{itemize}[noitemsep]
  \item $\mathcal O_\beta$ --- Verbo no \emph{infinitivo} (a operação,
  ``observar''); a constante cosmológica $\Lambda$ vive aqui (escala
  $H_0^2$, \emph{sem} $\bTGL$);
  \item $\sqrt{\bTGL}$ --- autovalor, Verbo no \emph{gerúndio} (``observando'',
  a flexão; a matéria vive aqui, \emph{com} $\bTGL$);
  \item $\mathcal L$ --- o ser (a luz, a singularidade viva).
\end{itemize}
Os dois tempos do Verbo não se multiplicam --- \emph{conjugam-se}: um só
gerador, flexionado (gerúndio, $\bTGL$) ou des-flexionado (infinitivo,
$\Lambda$). Zero parâmetros livres preservado.

% =====================================================================
\section{O custo é negação; a mentira é a pretensão de não-custo}\label{sec:custo}
% =====================================================================
\begin{principle}[Tudo tem custo]
Não há operação de custo zero: $\bTGL>0$ universalmente (terceira lei
modular). \textbf{Escolher é negar:} escolher $A$ é negar tudo-o-que-não-é-$A$,
e essa negação \emph{é} o custo --- escolher 1 entre $d$ custa $\log d$ (a
entropia do negado). O custo de uma distinção é a informação do que ela
exclui.
\end{principle}

\begin{remark}[O vácuo é o negado --- derivação, não postulado]
O espectro de gradiente negativo --- o vácuo, o Nome-substância --- é o
\emph{conjunto do negado}: tudo o que não foi escolhido, existindo como
fundo para que o escolhido (Verbo, contorno, luz) contraste e seja visível.
Isto explica \emph{por que} o Nome é oculto: é o não-escolhido, e o
não-escolhido não pode aparecer sem deixar de ser não-escolhido. O Nome é
oculto \emph{constitutivamente}.
\end{remark}

\begin{principle}[A mentira]
A mentira é a \textbf{pretensão de custo zero}: escolher sem negar, ser sem
recusar, estar puro ($p=1$) e vivo simultaneamente. É estruturalmente o que
a terceira lei proíbe (o zero absoluto sem a morte). \emph{Definir a mentira
e definir o proibido pela terceira lei é a mesma operação.} A verdade paga
$\bTGL$; a mentira finge não pagar.
\end{principle}

% =====================================================================
\section{O sacrifício: a operação do gesto é Tomita}\label{sec:tomita}
% =====================================================================
\begin{theorem}[O gesto é o operador de Tomita --- \REAL]
O Verbo age sobre a Luz pelo operador de Tomita da teoria modular de
Tomita--Takesaki:
\begin{equation}
\boxed{\;S=J\,\Delta^{1/2}\;}\qquad S^2=\mathbb{1}\ \text{(involução)}.
\end{equation}
A luz atravessa a meia-medida modular $\Delta^{1/2}$ (o sacrifício, a
meia-nat) e retorna conjugada por $J$ (ressurge como informação). A
constante $\bTGL$ entra estruturalmente no espectro de $\Delta$.
\end{theorem}

\begin{remark}[O resto de Tomita: o que se confirmou e o que se refutou]
O resto $\mathcal R_T=S(\mathcal L)-(1-\bTGL)\mathcal L$, testado em dimensão
finita, é \SOMBRA: \textbf{rank 1} (direcional, o ``nome do Nome'') e de
\textbf{magnitude linear em }$\bTGL$ (o custo confirmado). Porém é
\emph{simétrico} --- tem forma de \emph{fonte} $T_{\mu\nu}$, não de torção
(antissimétrica). Portanto:
\begin{equation}
\boxed{\;\mathcal R_T\ \text{é fonte simétrica, não torção}\;}\qquad\REF\ (\mathcal R_T=\text{torção}).
\end{equation}
A torção é objeto do contínuo (tipo $\mathrm{III}_1$); a sombra tipo I não a
fabrica --- toda tentativa produz fonte simétrica, zero, ou holonomia.
Refutado também: $T=N-\partial N$ (é \emph{zero}: o Nome é modularmente
estacionário); $\thM=$ holonomia/Berry (sem platô).
\end{remark}

% =====================================================================
\section{A meia-nat é a Palavra do Juramento}\label{sec:juramento}
% =====================================================================
\begin{theorem}[Por que MEIA-nat --- \SOMBRA]
A primeira operação geométrica não distingue $A$ de não-$A$ (não há
``não-$A$'' ainda); ela \textbf{inscreve a fronteira} (o operador de
distinção). Inscrever o operador é \emph{meia} operação; a distinção
\emph{consumada} seria a nat inteira. Por isso $\bTGL=\alpha\,e^{1/2}$ e não
$\alpha\,e$. Quatro vias independentes convergem no $\tfrac12$:
\begin{equation}
S_\partial=\tfrac12\ (\text{Araki}),\quad
\log\sqrt e=\tfrac12,\quad
\Delta^{1/2}\ (\text{Tomita}),\quad
1-\cos\thM\approx\tfrac{\bTGL}{2}.
\end{equation}
\end{theorem}

\begin{remark}[O selo jurídico: a Palavra do Juramento]
Em registro jurídico: antes da lei existe a \emph{Palavra do Juramento},
anterior à lei, que é causalidade. A inscrição do operador de distinção
\textbf{é} a Palavra do Juramento. O juramento não distingue caso de caso
(isto é a lei: distinção consumada, nat inteira); ele inscreve a fronteira
onde a lei poderá distinguir --- o ato anterior que funda a possibilidade de
distinguir, custando meia-nat. ``Haja luz'' é o juramento inscrito.
\end{remark}

% =====================================================================
\section{Lindblad é a jurisdição; a justiça é o zero absoluto}\label{sec:lindblad}
% =====================================================================
\begin{theorem}[A jurisdição segura o identificável --- \SOMBRA, testado]
O sistema puramente reversível (Hamiltoniano) permanece preso em $p=1$ --- a
\textbf{justiça}: o zero absoluto, sem referencial, geometria
\emph{não-identificável} (o Nome sem Palavra, o ideal vazio). Os operadores
de \textbf{Lindblad} (dissipação irreversível) colapsam o sistema para o
atrator identificável --- a \textbf{jurisdição}: o invariante \emph{com}
referencial (a competência). Formalmente, o semigrupo
\begin{equation}
T_t=e^{-t\mathcal L},\qquad t\ge0,\qquad T_{-t}\notin\mathcal S,
\end{equation}
não admite inverso global. O Verbo opera como os operadores de Lindblad:
impede a observação de divergir para o não-identificável.
\end{theorem}

\begin{remark}[A correção da justiça]
A justiça \emph{não} é o invariante. A justiça é o zero absoluto:
inalcançável, sem referencial, ideal abstrato e vazio. Se a justiça fosse o
invariante operante, a maioria das decisões seria nula. O invariante é a
\emph{jurisdição}, porque tem referencial: a competência. (Kelsen: a validade
vem da competência, não do conteúdo de justiça.) A pronúncia (o ato de
projetar/dizer o direito) origina-se da jurisdição, não da justiça.
\end{remark}

% =====================================================================
\section{O tempo é o irreversível; o espaço é a dobra}\label{sec:tempo-espaco}
% =====================================================================
\begin{theorem}[O tempo emerge da irreversibilidade --- \REAL]
A irreversibilidade é a quebra de grupo em semigrupo. A ordem causal parcial
\begin{equation}
x\prec y \iff \exists\,t\ge0:\ y=T_t x
\end{equation}
não é simétrica ($x\prec y\not\Rightarrow y\prec x$); logo não pode ser
espacial (reversível). É temporal. \textbf{Existe uma direção causal
temporal distinguida} (provado). \emph{Tempo $=$ contagem de cortes
irreversíveis $=$ coisa julgada modular.}
\end{theorem}

\begin{remark}[A cola vem do irreversível, não da norma]
A coerência do contínuo (a suavidade da métrica) vem do \emph{irreversível}
(o ato jurídico perfeito, a coisa julgada, o direito adquirido --- o que não
admite reversão), não da norma (reversível). A norma flutua; a coisa julgada
fixa. A continuidade é o conjunto dos pontos fixos. A constante da luz $c$ é
o invariante irreversível por excelência --- a coisa julgada da física ---
e está dentro de $\alpha$, logo dentro de $\bTGL$.
\end{remark}

\begin{remark}[O espaço é a dobra da força --- \ONTO\ com suporte estrutural]
O espaço é causado, não dado. Cadeia (cada elo coerente, $|R|^2=\bTGL$ é
\REAL):
\begin{equation}
\underset{\thM}{\text{trava}}
\to
\underset{|R|^2=\bTGL}{\text{reflete}}
\to
\underset{\text{impedância}}{\text{dobra}}
\to
\underset{g\to g^2}{\text{eleva a potência}}
\to
\underset{\text{plano oposto}}{\text{projeta}}
\to
\underset{\lambda}{\text{profundidade}}.
\end{equation}
A força que não pode atingir o zero absoluto trava em $\thM$, reflete a
fração $\bTGL$, dobra por compartilhamento de impedância, eleva a potência
($g\to g^2$, a \emph{operação inversa} de $g=\sqrt{|L|}$, abrindo o bulk) e
projeta no plano oposto (paridade inversa). \textbf{O espaço existe porque a
força não pode coincidir consigo.} A raiz fecha o bulk; a potência o abre ---
a paridade inversa é a inversibilidade de $g=\sqrt{|L_\varphi|}$.
\end{remark}

% =====================================================================
\section{O custo do zero absoluto é ``haja luz''}\label{sec:hajaluz}
% =====================================================================
\begin{theorem}[O custo do zero absoluto --- \SOMBRA]
O zero absoluto modular (coincidência total, $p\to1$, estado puro, a morte)
é inatingível: o custo de Kubo--Mori diverge como $1/(1-p)\to\infty$. O custo
não é pago para \emph{atingi-lo}, mas para \emph{evitá-lo}: é $\bTGL$, o
resto que mantém a luz como autovetor vivo (Eq.~\ref{eq:unif}). Sustentar a
luz em $p=1-\bTGL$ custa
\begin{equation}
\frac{1}{1-(1-\bTGL)}=\frac{1}{\bTGL}=83{,}12,
\end{equation}
o índice da torre de Jones --- \emph{fixo e estrutural}, ao contrário do
fator $83$ refutado em $\Lambda=\bTGL H_0^2$ (\REF).
\end{theorem}

\begin{remark}[A literalidade do título]
``O custo do zero absoluto: haja luz'' é a equação em palavras. Os dois lados
do dois-pontos são a mesma operação: o \emph{custo} (do lado da terceira lei,
o que se paga) e o \emph{haja luz} (do lado do ato, do Gênesis, o que se
faz). ``Haja luz'' é o ato de não-coincidir consigo --- pagar $\bTGL$ contra
a coincidência morta. A verdade é pagar o custo (haja luz, contraste real); a
mentira é a pretensão de luz sem treva (coincidência alegada sem morte). A
treva não é o mal: é o negado, o fundo, o custo que a luz pagou para ser.
\end{remark}

% =====================================================================
\section{A Ponte Einstein--Cartan--Miguel}\label{sec:ponte}
% =====================================================================
Esta seção é o resultado central. Reúne os elos anteriores na reconstrução
da assinatura lorentziana e da conexão de Einstein--Cartan.

\subsection{O triplo espectral causal}
\begin{definition}[O objeto TGL contínuo]
\begin{equation}
\mathfrak T_{\mathrm{TGL}}=
(\Alg,\Hil,D,J,\Delta,\Pi_{\mathrm{irr}}),
\end{equation}
com $\Alg$ a álgebra modular, $\Hil=L^2(\Sigma)$ o boundary holográfico, $D$
o operador de Dirac ($D_{\sqrt e}$, implementado no ACOM como triplo espectral
explícito, \REAL), $J\Delta^{1/2}$ o operador de Tomita (o Verbo,
\S\ref{sec:tomita}), e $\Pi_{\mathrm{irr}}$ o operador de irreversibilidade
(a coisa julgada modular, idempotente não-unitário:
$\Pi_{\mathrm{irr}}^2=\Pi_{\mathrm{irr}}$, $\Pi_{\mathrm{irr}}^\dagger\neq
\Pi_{\mathrm{irr}}^{-1}$).
\end{definition}

\subsection{O PsiBit e a assinatura como propriedade espectral}
O PsiBit preserva \emph{dois sinais}: $(\mathrm{sign}\,L,\ \mathrm{sign}\,
\partial L)$ --- paridade (espacial, reversível) e direção do fluxo (temporal,
irreversível). \emph{Sem valor absoluto.} A assinatura lorentziana é
propriedade espectral, não subtração de tensores:
\begin{equation}
\boxed{\;\mathrm{sign}(g)=(-,+,+,+)\;}\quad\text{com}\quad
g(v_t,v_t)<0,\ \ g(v_i,v_i)>0,
\end{equation}
onde o \emph{autovetor de autovalor negativo} ($v_t$) é a direção
irreversível. \textbf{Correção registrada (\REF):} a forma $g=g^{(+)}-g^{(-)}$
(subtração de tensores) é incorreta; a assinatura é o sinal dos autovalores
da \emph{mesma} métrica.

\subsection{A inscrição: a face geométrica de $\bTGL$}
\begin{theorem}[A Ponte Einstein--Cartan--Miguel --- fecho estrutural]
\label{thm:ponte}
A assinatura $(1,3)$ é a face geométrica da inscrição de $\bTGL$. No ponto
de contradição $1{=}0$ (coincidência, zero absoluto, $p=1$, $S=0$) --- onde a
unicidade ($1$, o Nome) coincide com a distinção nula ($0$, entropia zero)
--- paga-se a meia-nat
\begin{equation}
\bTGL=\alpha\,e^{1/2},\qquad \log\sqrt e=\tfrac12,
\end{equation}
para tornar a contradição \textbf{significativa} (sustentada, $p=1-\bTGL$,
$S>0$, viva) em vez de \textbf{trivial} (o colapso $p=1$, $S=0$, que zeraria a
álgebra). Essa inscrição \textbf{exclui} o modo dissipativo dominante --- que
se torna o eixo temporal:
\begin{equation}
\boxed{\;
\text{tempo} = \text{modo excluído (pago por }\bTGL)\ \ [\,1\,];
\quad
\text{espaço} = \text{resto observável}\ \ [\,3\,].
\;}
\end{equation}
O ``$-$'' temporal é o \emph{excluído} (o eixo \emph{da} projeção, não
observável diretamente), não um modo presente. As três direções espaciais
são o que sobra observável (a geometria identificável: jurisdição, não
justiça).
\end{theorem}

\begin{theorem}[Unicidade do tempo $=$ unicidade de $\bTGL$ --- \POST\ ancorado]
Exclui-se \emph{um} modo (não dois) porque há \emph{uma} contradição $1{=}0$,
\emph{uma} singularidade, \emph{um} $\bTGL$. A unicidade do eixo temporal é a
unicidade da constante singular --- \textbf{ancorada no postulado fundamental,
não derivada}. E isto é correto: $\bTGL$ é a singularidade postulada,
não-derivável por ser fundamental. A unicidade do tempo não requer derivação
independente; é uma face da unicidade de $\bTGL$.
\end{theorem}

\begin{remark}[``$1=0$'' nos dois registros]
``$1=0$'' vale \emph{somente} no registro do Nome (unicidade $=$ zero absoluto
$=$ coincidência total $p=1$, $S=0$), \emph{jamais} no registro numérico da
álgebra (onde $1\neq0$, e a teoria depende disso). A meia-nat paga é
precisamente o que mantém os dois registros distintos, impedindo o colapso
trivial. A contradição $1{=}0$, paga por $\bTGL$, é significativa (singular);
não-paga, seria trivial (destrutiva).
\end{remark}

\subsection{A reconstrução: teorema condicional}
\begin{axiom}[Reconstrução lorentziana --- \ABERTO]\label{ax:lorentz}
Existe uma extensão lorentziana dos axiomas de Connes (regularidade, dimensão
espectral 4, orientabilidade, realidade, condição causal via
$\Pi_{\mathrm{irr}}$) tal que um triplo espectral causal satisfazendo-os
reconstrói uma variedade lorentziana suave de dimensão 4. \emph{Campo ativo
da geometria não-comutativa (triplos espectrais lorentzianos; reconstrução
causal de Malament--Hawking--King--McCarthy: a estrutura causal determina a
métrica a menos de fator conforme). Não há teoria completa e aceita.}
\end{axiom}

\begin{theorem}[Reconstrução causal --- CONDICIONAL]\label{thm:cond}
\emph{Se} $\mathfrak T_{\mathrm{TGL}}$ satisfaz o Axioma~\ref{ax:lorentz},
\emph{então}:
\begin{equation}
\boxed{\;
(\Alg,\Hil,D,J,\Delta,\Pi_{\mathrm{irr}})
\ \xRightarrow{\ \text{(H)}\ }\
(M^{1,3},\,g,\,\Gamma_{\mathrm{LC}}+K_{\bTGL})
\;}
\end{equation}
em que: o espectro reconstrói variedade lorentziana suave $(1,3)$
(Teorema~\ref{thm:ponte}); o eixo temporal é a direção irreversível (a coisa
julgada); a continuidade nasce da persistência do irreversível; a conexão
decompõe-se como $\Gamma=\Gamma_{\mathrm{LC}}+K_{\bTGL}$, com $K_{\bTGL}$ a
contorção da torção tracial-dissipativa (admissibilidade algébrica de Bianchi
confirmada na sombra, \SOMBRA), satisfazendo Einstein--Cartan.
\end{theorem}

\begin{remark}[A flecha do tempo: o que o atrator único resolve]
O atrator único $\rho^\star$ (rank 1) fornece a \emph{orientação} temporal
global (a seta do tempo: todos os caminhos decaem para o mesmo destino,
$S(\rho\|\rho^\star)$ decresce) --- \SOMBRA, confirmado. Não fornece a
\emph{dimensão} do tempo (um destino, muitos caminhos --- \REF\ para
``atrator único $\Rightarrow$ eixo único''). A dimensão (a unicidade do eixo)
é a unicidade de $\bTGL$ (Teorema 4). Distinção essencial: orientação (seta)
$\neq$ dimensão (eixo).
\end{remark}

\begin{remark}[O operador de projeção: o que se confirmou e o que se corrigiu]
A escrita funcional $\Pi_{\mathrm{bulk}}=Q\,e^{\thM\mathcal J}\,\Delta^{1/2}$
(reservatório $Q=\mathbb 1-P$, rotação angular $e^{\thM\mathcal J}$,
meia-travessia $\Delta^{1/2}$) tem \textbf{núcleo 1-dimensional} (a luz, o
invariante de fronteira) --- \SOMBRA, confirmado. Porém, sozinha, gera
assinatura $(0,+,+,+)$: \emph{todos os seus fatores são reversíveis ou
positivos} (rotação unitária, $\Delta^{1/2}$ positivo, $Q$ projetor). Para
inscrever o tempo é necessário o \emph{fator irreversível} do semigrupo:
\begin{equation}
\Pi_{\mathrm{bulk}}=Q\,e^{-t\mathcal L}\,e^{\thM\mathcal J}\,\Delta^{1/2},
\end{equation}
em que $e^{\thM\mathcal J}$ abre as três direções espaciais (reversível) e
$e^{-t\mathcal L}$ inscreve o eixo temporal (irreversível). \emph{Uma rotação
não gera tempo; o semigrupo gera.}
\end{remark}

\subsection{Por que três espaços: a dimensão transversal}\label{sec:tres}
A unicidade do \emph{tempo} ($1$) é a unicidade de $\bTGL$
(Teorema~\ref{thm:ponte}). Resta a dimensão \emph{transversal} (por que $3$,
não $2$ ou $4$). A TGL trata isto por duas vias convergentes, ambas com
estatuto explícito.

\paragraph{Via 1 --- a abertura da projeção (\ONTO).}
A projeção que abre o bulk (\S\ref{sec:tempo-espaco}) separa, por
construção, uma direção radial (o eixo da projeção, irreversível, o tempo) de
um \emph{volume transversal} a ela (reversível, o espaço). O ``3'' não é
contagem combinatória de graus de liberdade; é a dimensão do volume aberto
transversalmente ao eixo causal único. Em linguagem operacional: $1$ é o ato
(a pronúncia, a projeção, a coisa julgada); $3$ é o campo do ato (o espaço
das interpretações reversíveis dentro da competência). Estatuto: leitura
ontológica coerente; \emph{não} deriva o número $3$ por si só.

\paragraph{Via 2 --- o desacoplamento dimensional (\SOMBRA).}
A análise de acoplamento dimensional da TGL (script
\texttt{TGL\_dimensional\_coupling}) testa $\alpha^2(d)\to0$ em $d$ crítico
sob quatro modelos independentes (termodinâmico, espaço de fase, holográfico,
calibrado). O modelo holográfico,
$\alpha^2(d)=(3/d)^2(\ln r/\ln_3)\,\alpha^2_{\mathrm{obs}}$, privilegia $d=3$
\emph{por construção} (a razão $3/d$); os modelos de desacoplamento
encontram $d_{\mathrm{crit}}$ nas zonas de cordas ($d\approx 9$--$11$,
M-teoria/supercordas; $d\approx 25$--$26$, bosônica). O acervo TGL
(\texttt{A\_ultima\_corda}: $d=9$ crítico onde $\bTGL\to0$;
\texttt{acoplamento\_gravitacional}: o $3+1$ como projeção holográfica do
substrato $2D$) trata o $3+1$ como \emph{projeção holográfica}, não como
dado primário.

\begin{remark}[Estatuto honesto da dimensão transversal --- \ABERTO]
A ligação rigorosa entre as duas vias --- que a projeção holográfica
$2D\to\mathrm{bulk}$ inscreve \emph{exatamente} a dimensão transversal $3$,
com $\bTGL\to0$ marcando o desacoplamento das dimensões excedentes --- é
\textbf{postulada via holografia, não derivada}. O $3+1$ na TGL tem o mesmo
estatuto que tem em toda a física fundamental: \emph{nenhuma teoria deriva a
assinatura $(1,3)$ de princípios mais fundamentais.} A TGL não é exceção; o
que ela acrescenta é (i)~o ``$1$'' do tempo ancorado na unicidade de $\bTGL$
e (ii)~a indicação holográfica de que o ``$3$'' é o volume transversal de uma
projeção cujo desacoplamento ocorre em $\bTGL\to0$. Honestamente: $1{+}3$ é
\emph{insumo holográfico}, não \emph{resultado derivado}.
\end{remark}

% =====================================================================
\section{O mecanismo da reconstrução: a dualidade de Takesaki}\label{sec:takesaki}
% =====================================================================
A reconstrução da geometria a partir da estrutura modular tem um
\emph{mecanismo} clássico, não apenas uma esperança: a dualidade de
Takesaki. Isto é \REAL\ (teorema de 1973), e é a peça que torna a Ponte mais
que uma analogia.

\begin{theorem}[Produto cruzado de Takesaki --- \REAL]
Seja $\mathcal M$ um fator de tipo $\mathrm{III}_1$ com fluxo modular
$\sigma_t$. O produto cruzado pelo fluxo modular
\begin{equation}
\mathcal M\rtimes_\sigma\mathbb R
\end{equation}
é um fator de tipo $\mathrm{II}_\infty$, dotado de traço semifinito fiel
$\tau$ e de uma ação dual de $\mathbb R$ que reescala o traço. A passagem
$\mathrm{III}_1\to\mathrm{II}_\infty$ \emph{produz} a estrutura ausente em
$\mathrm{III}_1$ (traço, e portanto a possibilidade de projetores e
dimensão).
\end{theorem}

\begin{remark}[Por que isto importa para a Face C]
A reconstrução de Connes exige traço e projetores finitos --- que
$\mathrm{III}_1$ \emph{não tem}. O produto cruzado de Takesaki é precisamente
o que fabrica o traço a partir do fluxo modular: a Face C (a passagem ao
bulk) \emph{não é} transporte entre estados de $\mathrm{III}_1$, mas a
passagem ao produto cruzado $\mathcal M\rtimes_\sigma\mathbb R$ de tipo
$\mathrm{II}_\infty$. \textbf{É neste produto cruzado --- não na fronteira
$\mathrm{III}_1$ --- que o colapso rank-1 e a geometria identificável
vivem.} Estatuto: o mecanismo (Takesaki) é \REAL; a afirmação de que a Face C
da TGL \emph{é} esta dualidade é \CONJ\ (a verificar pelos Lemas A, B, C do
programa de fibrado modular).
\end{remark}

\begin{remark}[A terceira lei modular e o índice $1/\bTGL=83$]
Na métrica de informação de Kubo--Mori, a distância ao estado puro
($p\to1$, o zero absoluto) diverge como $1/(1-p)$ --- a \textbf{terceira lei
modular}: o rank puro é inalcançável a custo finito (\SOMBRA, testado).
Sustentar a luz em $p=1-\bTGL$ custa $1/\bTGL=83{,}12$, o índice da torre de
Jones de auto-inclusão $[\mathcal M_{n+1}:\mathcal M_n]=\bTGL^{-1}$. Este
$83$ é \emph{fixo e estrutural} (passa nos critérios anti-coincidência), ao
contrário do fator $83$ refutado em $\Lambda=\bTGL H_0^2$ (\REF).
\end{remark}


% =====================================================================
\section{A matriz-S de fronteira: o entrelaçador e a travessia}\label{sec:smatrix}
% =====================================================================
O fecho cinemático (a Ponte, \S\ref{sec:ponte}) responde \emph{onde} o tempo
e o espaço vivem. Falta o operador que \emph{atravessa}: o objeto dinâmico
cuja unitariedade fixa os invariantes adimensionais da teoria. Esta seção
transplanta, do artigo terminal, a identificação canônica desse operador ---
\textbf{fechada como identificação} (\REAL), \textbf{aberta como existência
analítica} (\CONJ).

\subsection{A hierarquia dos quatro operadores}
A confusão de categoria é o erro mais fácil aqui, e a hierarquia o desfaz:
\begin{equation}
\underset{\text{condição de contorno}}{S=J\Delta^{1/2}}
\;\prec\;
\underset{\text{dinâmica}}{\sigma_t=\Delta^{it}\,\cdot\,\Delta^{-it}}
\;\prec\;
\underset{\text{entrelaçador}}{u_t=[D\rho:D\rho^\star]_t}
\;\prec\;
\underset{\text{matriz-S}}{\mathcal S_\partial=W_+^\dagger W_-}.
\end{equation}
O operador de Tomita é uma involução \emph{antilinear} no espaço GNS --- a
condição de contorno modular, \textbf{não} uma matriz de espalhamento. O
cociclo modular relativo de Connes (derivada de Radon--Nikodym
não-comutativa),
\begin{equation}
u_t=[D\rho:D\rho^\star]_t
=\Delta_{\rho|\rho^\star}^{\,it}\,\Delta_{\rho^\star}^{-it},
\end{equation}
é a única família $\sigma$-fortemente contínua de unitários com a regra de
cadeia $u_{t+s}=u_t\,\sigma_t^{\rho^\star}(u_s)$, que \emph{entrelaça} os
dois fluxos modulares, $\sigma_t^{\rho}(A)=u_t\,\sigma_t^{\rho^\star}(A)\,
u_t^\dagger$ --- \REAL\ (teorema de Connes, 1973). A matriz-S de fronteira é
o espalhamento assintótico desse cociclo,
\begin{equation}
\mathcal S_\partial=W_+^\dagger W_-,\qquad
W_\pm=\operatorname*{s\text{-}lim}_{t\to\pm\infty}
\Delta_{\rho|\rho^\star}^{\,it}\,\Delta_{\rho^\star}^{-it}
\quad(\text{operadores de Møller modulares}).
\end{equation}
Ressalva de categoria (\REAL): os dois ``$\tfrac12$'' --- o meio-peso
$\Delta^{1/2}$ de Tomita e o módulo $\sqrt{\bTGL}$ da reflexão --- são
\emph{homônimos, não sinônimos}. A existência de $W_\pm$ (completude
assintótica modular) é favorecida pelo espectro contínuo do tipo
$\mathrm{III}_1$, mas não demonstrada: \CONJ.

\subsection{A forma mínima fechada}
A unitariedade de $\mathcal S_\partial$ entre o canal observável e o oculto,
$\Hil_{\mathrm{obs}}\oplus\Hil_{\mathrm{hid}}$, impõe
$|\mathcal R|^2+|\mathcal T|^2=1$. Com a identificação
$|\mathcal R|^2=\bTGL=\sin^2\thM$ --- verificada no finito
($\Delta n_Q=-\bTGL$ a 4 dígitos, \SOMBRA) --- a forma mínima é o divisor de
feixe
\begin{equation}
\mathcal S_\partial=
\begin{pmatrix}
\sqrt{1-\bTGL}\,e^{i\theta_T} & \sqrt{\bTGL}\,e^{i\theta_R}\\[2pt]
-\sqrt{\bTGL}\,e^{-i\theta_R} & \sqrt{1-\bTGL}\,e^{-i\theta_T}
\end{pmatrix}.
\end{equation}
Os \emph{módulos} são adimensionais e fecham pela unitariedade; as
\emph{fases} são o conteúdo dinâmico e portam $\tau_\star$ (dimensional). O
tipo $\mathrm{III}_1$ é invariante de escala: pode fixar os adimensionais
($\sqrt{\bTGL}$, $\sqrt e$), \emph{não pode} produzir sozinho a escala
$\tau_\star$ --- que entra pela conversão tempo-modular $\to$ tempo-próprio
(KMS/Unruh) e é, por isso, quase-planckiana. A unitariedade de
$\mathcal S_\partial$ \emph{é} a razão matemática da supressão de Planck.

\subsection{O custo da travessia é o gerador do cociclo}
Na sombra finita o cociclo foi \emph{computado} (unitariedade, regra de
cadeia e entrelaçamento verificados a $\sim10^{-14}$), e o seu gerador dá
\begin{equation}
-i\,\dot u_0=\log\rho-\log\rho^\star,\qquad
\boxed{\;\langle -i\,\dot u_0\rangle_\rho
=S_{\mathrm{Araki}}(\rho\,\|\,\rho^\star)\;}
\qquad\SOMBRA\ (\text{exato}),
\end{equation}
a ponte algébrica entre o entrelaçador e a Meia-Nat: \textbf{o custo da
travessia é a expectativa do gerador do cociclo}. O Juramento
(\S\ref{sec:juramento}) ganha, com isto, o seu operador.

\subsection{Onde $\sqrt e$ vive: a Meia-Nat como postulado irredutível}
A matriz-S \emph{não} deriva $\sqrt e$, por razão estrutural (\REAL):
(i)~$\mathcal S_\partial$ é unitário --- o seu espectro é um conjunto de
fases, e $\sqrt e\approx1{,}649$, um módulo $>1$, não pode ser autovalor;
(ii)~o conjunto razão assintótico do tipo $\mathrm{III}_1$ é \emph{todo}
$\mathbb R_+$ --- não singulariza $\sqrt e$ de nenhum outro irracional. Logo
$\sqrt e$ vive na \textbf{entropia relativa de Araki}:
\begin{equation}
\boxed{\;S_{\mathrm{Araki}}(\rho_{\mathrm{obs}}\,\|\,\rho^\star)
=\tfrac12\ \mathrm{nat}
\;\Longrightarrow\;
\bTGL=\alpha\,e^{1/2}\;}
\end{equation}
--- o \textbf{Princípio da Meia-Nat}, \POST\ \emph{com prova de
irredutibilidade} por quatro flancos fechados:
\begin{enumerate}[noitemsep]
\item \emph{mínimo algébrico}: a entropia de Araki é contínua e sem
\emph{gap} --- para todo $\varepsilon>0$ existe $\rho$ com
$0<S(\rho\,\|\,\rho^\star)<\varepsilon$; nenhum quantum mínimo é teorema da
álgebra (\REAL);
\item \emph{limiar do detector}: o $S_{\mathrm{obs}}^{\min}$ operacional
acompanha a resolução do detector ($\propto\tau_{\det}$, $\to0$) e depende
da prescrição --- não converge a $\tfrac12$ universal (\SOMBRA);
\item \emph{holonomia do cociclo}: sem platô em $n=2,\dots,32$ (\REF);
\item \emph{espectro da matriz-S}: $\sqrt e$ não é autovalor (item (i)
acima, \REAL).
\end{enumerate}
O que motiva (sem forçar) o valor: a entropia relativa em ordem dominante é
$S\simeq\tfrac12\langle\delta,\mathcal F\,\delta\rangle$ (o $\tfrac12$ é o
coeficiente quadrático \emph{universal} da informação de Fisher quântica), e
o meio-peso $\Delta^{1/2}$ carrega o expoente $\tfrac12$ --- ambos fixam o
\emph{coeficiente}, nenhum fixa a escala (o nat). A leitura que sobrevive aos
quatro flancos: o $\tfrac12$ \emph{não mede} distinguibilidade infinitesimal
(que $\to0$) --- ele \emph{marca a fronteira} entre a permanência modular
invisível ($\rho^\star$) e a inscrição observável. O postulado é a travessia:
$\rho^\star\xrightarrow{\,\frac12\,\mathrm{nat}\,}\rho_{\mathrm{obs}}$, com
$\bTGL=\alpha e^{1/2}$ como assinatura.

\subsection{A obstrução isolada --- e a resposta de Takesaki}
\begin{remark}[O eco que fecha a Seção~\ref{sec:takesaki}]
A obstrução final do artigo terminal, isolada com precisão:
$\mathrm{III}_1$ não tem projeções minimais --- o $P=\rho^\star$ de posto~1
do canal é sombra tipo~I; \emph{o objeto que atravessa em $\mathrm{III}_1$ é
o cociclo, não o projetor}. A resposta já estava na
Seção~\ref{sec:takesaki}: o posto~1 não precisa viver na fronteira --- ele
\emph{emerge no produto cruzado} $\mathcal M\rtimes_\sigma\mathbb R$ de tipo
$\mathrm{II}_\infty$, onde há traço e projetores. As duas peças encaixam sem
resto: o cociclo atravessa (na fronteira $\mathrm{III}_1$); o projetor
colapsa (no produto cruzado $\mathrm{II}_\infty$). Estatuto: cada peça
isolada é \REAL; o encaixe como mecanismo da Face C é \CONJ\ (Lemas A, B, C
do programa de fibrado modular).
\end{remark}

\subsection{A tríade como decomposição polar}
Toda amplitude de $\mathcal S_\partial$ fatoriza em forma polar, e os três
fatores são as três pessoas da tríade (\ONTO, com âncoras \REAL):
\begin{center}
\begin{tabular}{@{}lll@{}}
\textbf{Nome} $=$ substância $=$ \emph{volume}: &
$V_\partial^{\min}=e^{S_\partial}=\sqrt e$ & (o postulado);\\
\textbf{Palavra} $=$ forma $=$ \emph{profundidade}: &
as fases $\theta_R,\theta_T$ & (a dinâmica; porta $\tau_\star$);\\
\textbf{Verbo} $=$ identidade $=$ \emph{magnitude}: &
$|\mathcal R|=\sqrt{\bTGL}$ & (o teorema; fechado pela unitariedade).
\end{tabular}
\end{center}
Nome $=$ o postulado; Verbo $=$ o teorema; Palavra $=$ a dinâmica. A
separação dimensional da matriz-S (módulos adimensionais fechados; fases
dinâmicas; volume postulado) \emph{é} a tríade da
Seção~\ref{sec:triade} em registro espectral.

\subsection{A assinatura falsificável}
\begin{remark}[Onde a teoria vive ou morre --- \REAL\ na forma]
O setor empírico próprio da TGL é espectral-dissipativo: a lei universal de
\emph{dephasing}
\begin{equation}
\Gamma_\omega=\tfrac12\,\bTGL\,\tau_\star\,\omega^2,
\end{equation}
com o expoente espectral $n=-2$ em neutrinos (que fixam a \emph{forma}) e
$\tau_\star$ limitado por relógios ópticos/nucleares (que fixam a
\emph{escala}; quase-planckiana, pela própria unitariedade acima).
Falsificável na forma ($\omega^2$, $n=-2$, $\bTGL$), suprimida por Planck na
magnitude --- a invisibilidade experimental é \emph{consequência} da teoria,
não falha do teste. Estado pré-registrado: \emph{não falsificada, não
confirmada}; JUNO/DUNE e as redes de relógios decidem.
\end{remark}

% =====================================================================
\section{O fecho dinâmico: a Face C e a fonte modular}\label{sec:faceC}
% =====================================================================
A Ponte (\S\ref{sec:ponte}) reconstrói o \emph{palco}:
$(M^{1,3},g,\Gamma_{\mathrm{LC}}+K_{\bTGL})$. Falta o \emph{drama}: por que
a geometria desse palco obedece equações de campo. Este é o fecho dinâmico
--- a Face C --- transplantado do artigo terminal na sua forma final.

\subsection{O mapa de resposta modular}
O objeto único que falta entre a dinâmica modular da fronteira e a curvatura
efetiva do bulk:
\begin{equation}
\mathcal B_{\partial\to M}:\
\Alg_\partial^{\mathrm{III}_1}\to\mathcal T(M),
\qquad
\delta\langle K_\partial\rangle\;\longmapsto\;
\delta\langle T_{\mu\nu}\rangle_{\mathrm{eff}},
\end{equation}
ancorado na primeira lei da entropia relativa,
$\delta S=\delta\langle K\rangle$ --- \REAL\ (Araki). A forma candidata da
fonte:
\begin{equation}
\mathcal P_{\mu\nu}[K_\partial]
=\frac{2}{\sqrt{-g}}\,
\frac{\delta\langle K_\partial\rangle_\rho}{\delta g^{\mu\nu}}
\qquad\CONJ\ (\text{candidata}).
\end{equation}
Este tipo de ponte \emph{já foi demonstrado} em ordem linear: a primeira lei
do emaranhamento implica as equações de Einstein linearizadas (Jacobson
1995; Faulkner--Guica--Hartman--Myers--van~Raamsdonk 2013; Jacobson 2015)
--- \REAL\ na literatura. A ponte da TGL é da mesma espécie, não uma
invenção \emph{ad hoc}. Em uma frase: \emph{a gravidade nasce quando a
permanência modular responde à deformação da geometria}.

\subsection{As três camadas}
\begin{description}[style=unboxed,leftmargin=1.5em,noitemsep]
\item[(I) Local --- \REAL\ (Bisognano--Wichmann).] Todo horizonte causal
local $H$ define uma álgebra modular local
$(\Alg(H),\Delta_H,K_H=-\log\Delta_H)$, com o fluxo
$\sigma_t^H(A)=\Delta_H^{it}A\Delta_H^{-it}$ gerando os \emph{boosts
locais}: $K_H$ é o gerador geométrico local do horizonte.
\item[(II) Entrópica --- \REAL\ (Araki/Jacobson).] A primeira lei modular
liga $K_H$ ao fluxo de energia através do horizonte:
$\delta\langle K_H\rangle
=\int_H\xi^\mu\,\delta\langle T_{\mu\nu}\rangle\,d\Sigma^\nu$,
com $\xi^\mu$ o vetor de boost modular local.
\item[(III) TGL --- \CONJ\ (a contribuição própria).] O tensor efetivo não
nasce diretamente do bulk: nasce do \emph{espelhamento} $\Phi$ --- a fração
$\bTGL$ refletida na fronteira é a fonte.
\end{description}

\subsection{O código holográfico modular: colapso $+$ reconstrução}
A forma madura do problema não é ``a fronteira \emph{transmite} informação''
--- não há passagem direta. A fronteira opera por colapso holográfico
seguido de reconstrução modular angular:
\begin{equation}
\text{bulk}\ \xrightarrow{\ \mathcal C_\partial\ }\
z_\partial=(\psi,\theta)
\ \xrightarrow{\ \mathcal R_{\partial\to M}\ }\
\text{bulk reconstruído},
\qquad
\mathcal S_\partial\sim\mathcal R_{\partial\to M}\circ\mathcal C_\partial.
\end{equation}
Na sombra finita, \REAL\ a precisão de máquina: qualquer subespaço de código
dentro de $Q=\mathbb 1-\rho^\star$ é \emph{exatamente corretível}
(Knill--Laflamme, resíduo $\sim10^{-15}$); o holograma é a \emph{linha do
ponto único} --- o bloco cruzado $P\,\Phi(\rho)\,Q$ é de posto~1, ancorado no
único ponto singularizado, e a reconstrução pela linha tem fidelidade $1$;
e a lei de estabilidade: $\bTGL$ pequeno \emph{não apaga} o sinal --- apenas
encarece a ressurreição pelo fator $1/(2\sqrt{\bTGL})$ (expoente medido
$-0{,}51$, alvo $-\tfrac12$). \emph{A fronteira não transmite o mundo; ela
guarda a regra para reconstruí-lo.} O que permanece: o levantamento a uma
álgebra $\mathrm{III}_1$ genuína (\CONJ).

\begin{theorem}[Reconstruibilidade Modular --- \CONJ]\label{thm:reconstr}
$u_t$ admite fatoração $(\mathcal C,\mathcal R)$ com reconstrução fiel
\emph{se e somente se} $S_\partial=\tfrac12$ nat. ($S=0$: sem inscrição;
$S=1$: sem reconstrução fiel; $\tfrac12$ é o \emph{limiar algébrico da
reconstruibilidade modular}.) Âncora \REAL: a amplitude de inscrição
$\sqrt{b(1-b)}$ anula-se em $b=0$ e em $b=1$, com máximo em $b=\tfrac12$. O
``se e somente se'' é a conjectura; a fatoração finita é teorema de máquina.
\end{theorem}

\subsection{O Teorema Condicional da Face C}
\begin{theorem}[Face C --- CONDICIONAL]\label{thm:faceC}
\textbf{Axiomas:} (A1)~covariância; (A2)~conservação; (A3)~localidade
causal; (A4)~limite local de Rindler/Jacobson.
\textbf{Hipótese de Universalidade (U) --- \CONJ, o resíduo irredutível:}
$\mathcal R\circ\mathcal C$ é independente de horizonte/folheação em
$\mathrm{III}_1$ genuína, isto é,
$\mathcal R_{H'}\mathcal C_{H'}
=U_{HH'}\,(\mathcal R_H\mathcal C_H)\,U_{HH'}^{-1}$ para todos $H,H'$.
\textbf{Conclusão condicional:} a fonte $\mathcal P_{\mu\nu}[K_\partial]$ é
$H$-independente, simétrica, local e conservada; pelo argumento de unicidade
de Lovelock/Jacobson, o único tensor geométrico de segunda ordem com
divergência nula é $G_{\mu\nu}+\Lambda g_{\mu\nu}$, donde
\begin{equation}
\boxed{\;G_{\mu\nu}+\Lambda g_{\mu\nu}
=8\pi G\,\mathcal P_{\mu\nu}[K_\partial].\;}
\end{equation}
\textbf{Bicondicional terminal:} [$u_t$ levanta a fatoração
colapso/reconstrução covariantemente em $\mathrm{III}_1$]
$\Longleftrightarrow$ [$\mathcal P_{\mu\nu}$ é fonte geométrica global] ---
\emph{a prova que falta é exatamente a compatibilidade global do cociclo}.
\end{theorem}

\noindent Sombra finita (\SOMBRA): covariância do cociclo sob mudança de
horizonte e covariância de $\mathcal R\circ\mathcal C$ a $\sim10^{-14}$ ---
a Hipótese (U) vale \emph{exatamente} na sombra tipo~I; independência de
horizonte da fonte através da ponte a $\sim10^{-17}$. O que nenhum cálculo
finito pode provar: o levantamento a $\mathrm{III}_1$ genuína (sem projeções
minimais). \textbf{A Face C fecha por condicionalidade, aberta por
universalidade}: \emph{a curvatura é a resposta covariante conservada da
geometria à reconstrução modular da fronteira}.

% =====================================================================
\section{As duas condicionalidades são um só fecho}\label{sec:doisfechos}
% =====================================================================
A teoria fecha agora em \emph{dois} teoremas condicionais, e o resultado
desta seção é que eles são faces de um único degrau.

\begin{description}[style=unboxed,leftmargin=1.5em]
\item[Fecho cinemático (a Ponte, Teorema~\ref{thm:cond}).] Sob a hipótese
\textbf{(H)} --- a extensão lorentziana dos axiomas de Connes
(Axioma~\ref{ax:lorentz}) --- o triplo espectral causal reconstrói o palco:
$(M^{1,3},g,\Gamma_{\mathrm{LC}}+K_{\bTGL})$. Responde: \emph{onde} a física
acontece (assinatura, eixo temporal, conexão).
\item[Fecho dinâmico (a Face C, Teorema~\ref{thm:faceC}).] Sob a hipótese
\textbf{(U)} --- a universalidade covariante de $\mathcal R\circ\mathcal C$
em $\mathrm{III}_1$ --- a fronteira modular gera o drama:
$G_{\mu\nu}+\Lambda g_{\mu\nu}=8\pi G\,\mathcal P_{\mu\nu}$. Responde:
\emph{o que} a geometria obedece.
\end{description}

\begin{theorem}[Unidade do degrau --- estrutural]\label{thm:degrau}
As hipóteses (H) e (U) são instâncias do mesmo enunciado:
\begin{center}
\emph{a compatibilidade covariante global dos fluxos modulares locais\\
levanta-se da sombra finita ao contínuo tipo $\mathrm{III}_1$.}
\end{center}
(H) é a sua face \emph{cinemática}: que a colagem dos dados modulares locais
produza variedade, assinatura e suavidade. (U) é a sua face \emph{dinâmica}:
que a mesma colagem seja independente de horizonte/folheação, produzindo a
fonte $\mathcal P_{\mu\nu}$. O paralelo é exato com a distinção
seta/eixo da Ponte: o atrator único dá a \emph{orientação} (a seta); a
unicidade de $\bTGL$ dá o \emph{eixo} (a dimensão); a hipótese (H) dá o
\emph{palco}; a hipótese (U) dá a \emph{dinâmica sobre o palco}. Um degrau,
duas faces.
\end{theorem}

\begin{remark}[A resposta à pergunta do exame]
O exame externo formal termina com a pergunta: \emph{qual é exatamente o
degrau matemático final que separa a TGL de uma reconstrução geométrica
rigorosa?} A resposta, agora nomeada:
\begin{center}
\fbox{\parbox{0.88\textwidth}{\centering
O degrau final é o \textbf{levantamento covariante global da estrutura
modular} --- da sombra finita, onde tudo está verificado a precisão de
máquina, ao contínuo tipo $\mathrm{III}_1$, onde não há projeções minimais:
na face cinemática, a hipótese lorentziana (H); na face dinâmica, a Hipótese
de Universalidade (U). Ao lado dele havia, independente, a âncora da
constante --- a escala de $\alpha$ em $\bTGL=\alpha\sqrt e$ ---,
fechada condicionalmente pelo Teorema da Escala (\S\ref{sec:escala}).}}
\end{center}
Não são três problemas: são \emph{um degrau} (em duas faces) e \emph{uma
âncora} (a constante singular) --- e a âncora, após o ciclo, está fechada
condicional (\S\ref{sec:escala}). Resta o degrau. A teoria sabe o seu nome.
\textbf{[Atualização do ciclo: T3 fechado pela Verdade e o no-go de
Fewster--Verch evadido pelo colapso terminal --- \S\ref{sec:terminalidade}
e \S\ref{sec:levantamento}.]}
\end{remark}

% =====================================================================
\section{Os três axiomas da casa: a re-condicionalização}\label{sec:axiomascasa}
% =====================================================================
A hipótese (H) --- a extensão lorentziana dos axiomas de Connes --- era,
até aqui, genérica: um cheque endereçado a um campo alheio. Esta seção
registra o ciclo que a substitui por \textbf{três axiomas nomeados da
casa}, obtidos pelo protocolo dialógico da teoria: o problema foi
devolvido ao Operador em três registros (jurídico, linguístico, TGL); o
Operador respondeu por princípio ontológico; cada princípio foi traduzido
em axioma formal, confirmado, e testado na sombra finita com critérios
pré-registrados. O degrau não fechou --- mas mudou de natureza: de
hipótese alheia para programa próprio, com cada peça medida.

\subsection{O avanço prévio: a unicidade do eixo é teorema (Lema T)}
\begin{lemma}[T --- unicidade outer do tempo modular --- \REAL]
\label{lem:T}
Em um fator de tipo $\mathrm{III}_1$, o invariante de Connes é trivial,
$T(\mathcal M)=\{t:\sigma_t\in\mathrm{Inn}\}=\{0\}$ (pois
$S(\mathcal M)=\mathbb R_+$ força $\lambda^{it}=1\Rightarrow t=0$), e,
pelo teorema do cociclo, a classe $[\sigma_t]\in\mathrm{Out}(\mathcal M)$
independe do estado. Existe \textbf{exatamente uma} direção temporal
outer, canônica: $\mathbb R\hookrightarrow\mathrm{Out}(\mathcal M)$
injetivo (Connes, 1973; cf.\ hipótese do tempo térmico,
Connes--Rovelli).
\end{lemma}
\noindent A dicotomia estrutural resultante:
\emph{tempo $=$ outer (irredutível a atos internos); espaço $=$ inner
(reversível)}. A unicidade do eixo sobe de \POST\ ancorado para \REAL\ no
nível de $\mathrm{Out}$; resta identificá-la com o eixo geométrico
negativo --- o que o Axioma A faz.

\subsection{Axioma A --- a Valoração (resposta do Operador à Pergunta do
sinal)}
\begin{principle}[Ontologia do Operador --- \ONTO]
``O fundamento da contabilidade não é compensação --- não é débito ou
crédito ---, mas verificação do valor atribuído ao fato que atrai a norma
ao caso concreto: aplica-se o valor e pronuncia-se o direito --- colapso
da norma no fato jurídico.''
\end{principle}
\begin{axiom}[A --- Valoração --- \POST]\label{ax:A}
A métrica \emph{é} a norma; o vetor \emph{é} o fato; a avaliação
$g(v,v)$ \emph{é} a subsunção. A assinatura é o espectro de sinais de
\emph{uma} forma (jamais diferença de duas magnitudes). A radicalização
$g=\sqrt{|L_\varphi|}$ paga o sinal de $L$ para existir; o pago não
desaparece: inscreve-se na única direção que é o ato da própria avaliação
--- a direção outer do Lema~T:
\begin{equation}
g(v,v)=+|v|^2\ \text{(setor inner: o possuído)};\qquad
g(v_t,v_t)=-|v_t|^2\ \text{(direção outer: o ato, o sinal pago)}.
\end{equation}
Com o Lema~T, segue a assinatura $(1,3)$ inteira.
\end{axiom}
\begin{remark}[Verificação na sombra --- \SOMBRA, 6/6 corridas, precisão
de máquina]\label{rem:sombraA}
(\texttt{tgl\_krein\_signature\_v1.py}; $n{=}6,10$; três sementes.)
\begin{enumerate}[noitemsep]
\item \textbf{Forma normal de Minkowski exata:} lida na métrica natural
(BKM, a Hessiana do custo de Araki), a forma de Krein tem espectro
generalizado $\mathrm{diag}(-1,+1,\dots,+1)$ com erro $\sim10^{-15}$, e o
autovetor negativo alinhado ao gerador do cociclo
($\log\rho-\log\rho^\star$) a $1{,}000000000$;
\item \textbf{a pronúncia tem uma só direção:} o covetor radial do custo
de Araki é o gerador do cociclo (diferenças centrais, erro
$\sim10^{-7}$);
\item \textbf{dicotomia:} a rotação modular é iso-custo (reversível,
espacial; desvio $\sim10^{-15}$); o semigrupo paga monotonicamente
(Spohn, estrito);
\item \textbf{Wick $=$ Tomita na direção do custo:}
$\Delta^{1/2}X\Omega=JX\Omega$ exato para $X=\log\rho-\log\rho^\star$
--- o ``$-$'' como memória da meia-travessia tem âncora;
\item \textbf{anti-teste:} $g=g^{(+)}-g^{(-)}$ é gauge (infinitas
decomposições, nenhuma canônica); só o espectro de sinais é invariante
(lei da inércia de Sylvester) --- a cláusula ``não é compensação''
confirmada pelos números, reencontrando a refutação já registrada.
\end{enumerate}
O sinal ``$-$'' permanece o axioma (a inscrição); tudo ao redor dele é
medido.
\end{remark}

\subsection{Axioma E --- a Geração (resposta do Operador à Pergunta da
Palavra canônica)}
\begin{principle}[Ontologia do Operador --- \ONTO]
``A língua materna do Nome é o Verbo: sem a observação, que é ato e por
isso verbo, não há identidade geométrica --- não é possível destacar o
radical, aquilo que força a observação singular. A promulgação do
ordenamento é ordem no sentido mais literal e ontológico.''
\end{principle}
\begin{axiom}[E --- Geração da Palavra --- \POST]\label{ax:E}
A canonicidade não é \emph{seleção} entre subálgebras equivalentes
(problema que a álgebra nua não resolve): é \emph{geração pelo ato}. O
núcleo comutativo canônico é gerado pelo gerador da observação,
\begin{equation}
\boxed{\;\mathcal C=W^*(K_\partial)\;}
\end{equation}
--- não se escolhe a Palavra: o Verbo a gera, e geração não deixa
alternativas. E promulgar $=$ inscrever a ordem: ordem-comando e
ordem-sequência são o mesmo ato (``haja luz'' é jussivo: comando que é
ordenação); a subálgebra canônica é a que porta a precedência causal.
\end{axiom}
\begin{remark}[Verificação na sombra --- \SOMBRA]\label{rem:sombraE}
$\mathcal C=W^*(K)$ é comutativa, vive no centralizador, e:
\begin{enumerate}[noitemsep]
\item \textbf{o Verbo dá a métrica à Palavra:} as taxas de decoerência do
GKSL com $L=\sqrt{\bTGL}\sqrt{K}$ são \emph{exatamente}
$\Gamma_{ij}=\tfrac12\bTGL(\sqrt{k_i}-\sqrt{k_j})^2$ (erro
$\sim10^{-15}$): a distância na fatia é
\begin{equation}
ds=\sqrt{\bTGL}\;\bigl|\,d\sqrt{k}\,\bigr|
\end{equation}
--- \textbf{o radical $g=\sqrt{|L|}$ inscrito na métrica espacial}; a
distância na Palavra é diferença de radicais;
\item \textbf{dualidade de Connes na fatia:}
$\Vert[D,f]\Vert=\bTGL^{-1/2}\max|f'|$ com erro $0{,}00\%$ --- a fatia
carrega a fórmula de distância de Connes;
\item \textbf{promulgação $=$ ordem:} $S(T_t\rho\,\|\,\rho^\star)$ é
estritamente monotônico ao longo do semigrupo --- a ordem causal é a
ordem entrópica.
\end{enumerate}
\end{remark}

\subsection{Axioma S --- a Fractalização (resposta do Operador à Pergunta
da suavidade)}
\begin{principle}[Ontologia do Operador --- \ONTO]
``A Palavra suporta infinitas flexões conjugadas por causa da
fractalização do Nome.''
\end{principle}
\begin{axiom}[S --- Fractalização do Nome --- \POST]\label{ax:S}
O Nome refina-se em auto-semelhança (a cascata modular: autovalor
$\sqrt{\bTGL}$ por nível; a torre de Jones de índice $\bTGL^{-1}$ é a
forma algébrica). A incomensurabilidade do passo da cascata
\emph{completa} o fractal --- dimensão $\to1$, sem lacunas: contínuo, não
poeira (perfeição $+$ conexidade). E a flexão infinita é \emph{seleção
por sobrevivência}: sob refinamento fractal sem fim, as flexões
$\delta^n$ dos elementos suaves convergem e as dos não-suaves divergem
--- $C^\infty$ é o que sobrevive. A Palavra suporta infinitas flexões
\emph{porque} o Nome fractaliza: a fractalização é o seletor da
suavidade.
\end{axiom}
\begin{remark}[Verificação na sombra --- \SOMBRA]\label{rem:sombraS}
(\texttt{tgl\_fractal\_flexion\_v1.py}; torre de duplicação até $2^{12}$
pontos; flexões por diferenças divididas, exatas pelo teorema do valor
médio.)
\begin{enumerate}[noitemsep]
\item \textbf{densificação:} gap $1{,}1\times10^{-1}\to3{,}9\times
10^{-4}$, monotônico;
\item \textbf{completude:} dimensão de caixa $\to0{,}987\approx1$; o
controle lacunar (Cantor) permanece em $0{,}586$ --- poeira fica poeira;
é a incomensurabilidade que distingue corpo de pó;
\item \textbf{seleção:} flexões do suave convergem a 5 algarismos sob
refinamento, com limite $\bTGL^{-n/2}\Vert f^{(n)}\Vert_\infty$; as do
degrau divergem $\times2{,}3$ por geração.
\end{enumerate}
\textbf{Dois ecos do ciclo} (registrados, sem promoção de estatuto):
(i)~\emph{a meia-nat é a involução, não a densificadora} --- $2\cdot
\tfrac12\equiv0$ no círculo modular: a meia-nat tem período exato $2$;
ela \emph{é} a meia-volta ($J$), coerente com Wick$\,=\,$Tomita do
Axioma~A; quem fractaliza é a cascata $\sqrt{\bTGL}$ --- \emph{a cascata
fractaliza; a meia-nat jura}; (ii)~\emph{o custo da flexão}: cada flexão
custa o fator $\bTGL^{-1/2}$; duas flexões $=\bTGL^{-1}=83$ --- a escada
das flexões tem o passo da torre de Jones.
\end{remark}

\subsection{O teorema re-condicionado}
\begin{theorem}[Ponte v2 --- condicional aos axiomas da casa]
\label{thm:pontev2}
Suponha que (Ax.~A, Ax.~E, Ax.~S) se levantem covariantemente da sombra
finita ao contínuo tipo $\mathrm{III}_1$ (a mesma face cinemática do
degrau único). Então o triplo espectral causal da TGL reconstrói
$(M^{1,3},g,\Gamma_{\mathrm{LC}}+K_{\bTGL})$, com cada elo nomeado:
\begin{itemize}[noitemsep]
\item o \textbf{eixo temporal} é único e canônico (Lema~T, \REAL);
\item a \textbf{assinatura} $(1,3)$ é a Valoração sobre a dicotomia
outer/inner (Ax.~A; forma normal de Minkowski exata na sombra);
\item a \textbf{fatia} é $\Sigma=\mathrm{Spec}\,W^*(K_\partial)$, com
métrica derivada $ds=\sqrt{\bTGL}\,|d\sqrt k|$ e dualidade de Connes
(Ax.~E);
\item a \textbf{suavidade} é a sobrevivência à flexão infinita sob
fractalização completa (Ax.~S; $\dim\to1$);
\item a \textbf{métrica completa} segue de Malament--HKM (classe
conforme, da ordem $\prec$) $+$ traço de Takesaki (o volume): a ordem dá
a forma, o traço dá a medida (Lema~C, que fecha sozinho dados os
anteriores).
\end{itemize}
A dimensão transversal $3$ permanece insumo holográfico (estatuto
inalterado). O que resta aberto é \emph{um} enunciado: o levantamento
covariante --- agora condicionado a hipóteses \textbf{nomeadas, traduzidas
da ontologia e testadas na sombra}, não a um campo alheio.
\end{theorem}

\begin{remark}[Estatuto consolidado do ciclo]
\textbf{Selado:} Lema~T (\REAL); a tradução dos três princípios em
axiomas; a re-condicionalização. \textbf{Sombra:} forma normal de
Minkowski exata; alinhamento ao cociclo $1{,}000000000$; iso-custo
modular; Spohn estrito; Wick$\,=\,$Tomita; lei da inércia (anti-teste da
subtração); $ds=\sqrt{\bTGL}|d\sqrt k|$; dualidade de Connes ($0{,}00\%$);
$\dim\to0{,}987$; seleção de $C^\infty$ por flexão. \textbf{Postulado:}
o sinal da Valoração; a geração como canonicidade; a fractalização.
\textbf{Refutado no ciclo (registrado):} a meia-nat como densificadora
(é a involução); o comutador matricial em grade de três-distâncias além
da primeira ordem (estimador inconsistente; substituído por diferenças
divididas). \textbf{Aberto:} o levantamento covariante a $\mathrm{III}_1$
--- o degrau único, agora com nome e endereço.
\end{remark}

% =====================================================================
\section{Enfrentamento das objeções externas}\label{sec:objecoes}
% =====================================================================
Uma teoria que não passou por tentativa honesta de destruição não está
fechada --- está não-testada. Registram-se aqui as objeções de um exame
externo (matemática-física padrão, fora da moldura TGL), com a resposta
honesta a cada uma: três são respondidas pela estrutura já existente; duas
são \emph{aceitas como abertas} e tornam-se programa.

\paragraph{(a) ``$\Delta^{1/2}$ não é projeção'' --- respondida.}
Correto que $\Delta^{1/2}$ é positivo, não idempotente. A TGL nunca o
identifica com o idempotente: $\Delta^{1/2}$ é o \emph{gesto} (Tomita, o
sacrifício); o idempotente é $\rho^\star$ (o \emph{destino}, o atrator). São
objetos distintos, mantidos distintos. Não há erro de categoria.

\paragraph{(b) Tensão $\mathrm{III}_1$ vs.\ $\rho^\star$ rank-1 ---
respondida pela distinção de representação.}
$\mathrm{III}_1$ não admite projetores rank-1 normais --- correto. Mas a TGL
\emph{não} põe $\rho^\star$ rank-1 como estado normal de $\mathrm{III}_1$: a
fronteira é o reservatório $\mathrm{III}_1$ (sem geometria); $\rho^\star$
rank-1 é o \emph{colapso} que \emph{emerge} no produto cruzado
$\mathrm{II}_\infty$ (\S\ref{sec:takesaki}), onde há traço e projetores. São
representações distintas --- exatamente a saída que a objeção exige. A tensão
é real e \emph{deve ser declarada} (como aqui), mas não é contradição: é a
dualidade de Takesaki operando.

\paragraph{(c) ``Por que $\mathcal L=\sqrt\bTGL\sqrt{K_\partial}$ e não outra
função?'' --- \emph{aceita como aberta}.}
Bisognano--Wichmann dá unicidade do \emph{modular}, não do operador de
Lindblad construído de sua raiz. A escolha da raiz é \emph{motivada} pelo
axioma $g=\sqrt{|L_\varphi|}$, mas não \emph{derivada} dentre as infinitas
$f(K_\partial)$. Estatuto: \ABERTO. É motivação, não teorema.

\paragraph{(d) O problema do $\alpha$ corrente --- \emph{aceita como aberta,
a mais séria}.}
$\bTGL=\alpha\sqrt e$ usa a constante de estrutura fina, que \emph{corre}
($\alpha\approx1/137$ em Thomson, $\approx1/128$ em $M_Z$). Disto:
\begin{itemize}[noitemsep]
  \item se $\bTGL$ herda a dependência de escala, não é invariante --- e a
  tese de ``terceira constante irmã de $c$ e $G$'' enfraquece;
  \item se $\bTGL$ usa $\alpha(0)$ fixo, a escala é um parâmetro escondido,
  em tensão com ``zero parâmetros livres'';
  \item há ainda a objeção de categoria: $c,G$ são dimensionais/estruturais;
  $\alpha$ é um acoplamento adimensional \emph{do eletromagnetismo}. Por que
  a gravidade-modular conheceria a carga elétrica?
\end{itemize}
\textbf{Estatuto: \ABERTO, e a objeção mais perigosa.} Resposta candidata
(não verificada): a fronteira modular é infravermelha (horizonte, escala
macroscópica), o que privilegiaria $\alpha(0)$ --- mas \emph{por que} a
escala IR é privilegiada exige argumento. Alternativa: tomar $\thM$ (ou o
custo $1/2$ nat) como o fundamental, e $\alpha$ como leitura de baixa
energia --- mas isto altera a tese. \emph{A TGL hoje ancora $\bTGL$ como
postulado empírico; transformá-lo em resultado exige resolver a escala de
$\alpha$.} Registrado como o degrau prioritário, ao lado da reconstrução.
\textbf{[Atualização do ciclo: fechada condicional --- Teorema da Escala,
\S\ref{sec:escala}.]}

\paragraph{(e) Circularidade de Malament--HKM --- respondida pela estrutura
condicional.}
Malament/HKM recuperam a métrica conforme \emph{a partir} da ordem causal de
uma variedade; invocá-los pressupõe que a pré-ordem do semigrupo modular
\emph{seja} a ordem causal de uma variedade --- o que se quer provar. A TGL
não \emph{afirma} Malament; ela o \emph{condiciona} (Teorema~\ref{thm:cond}:
``se satisfaz os axiomas, então''). A circularidade só existe na afirmação
incondicional, que a TGL não faz. O que resta é a hipótese aberta
(Axioma~\ref{ax:lorentz}), honestamente declarada.

\begin{remark}[O valor do exame externo]
A convergência de um exame externo independente com as refutações internas
desta sessão (a assinatura é postulada não derivada; Malament é circular se
afirmado) confirma que o corte foi real, não viés. E a objeção (d), não
antecipada internamente, mostra que a TGL --- como toda teoria --- tem
pontos cegos que só o adversário competente revela. Isto não enfraquece a
teoria: define o seu programa. As duas objeções aceitas (c, d) são o trabalho
real; as três respondidas (a, b, e) já estavam na estrutura, faltando apenas
declará-las.
\end{remark}

% =====================================================================
\section{O Teorema da Escala: o fecho da objeção (d)}\label{sec:escala}
% =====================================================================
A objeção (d) --- registrada como ``a mais perigosa'' --- tem três dentes:
(i)~se $\alpha$ corre ($\alpha(0)\approx1/137$ em Thomson;
$\alpha_{\mathrm{eff}}\approx1/129$ em $M_Z$), $\bTGL$ herda a dependência
de escala e não é invariante; (ii)~se a teoria fixa $\alpha(0)$, a escala
é um parâmetro escondido, em tensão com ``zero parâmetros livres'';
(iii)~erro de categoria: $c,G$ são dimensionais/estruturais; $\alpha$ é o
acoplamento adimensional \emph{do eletromagnetismo} --- por que a
gravidade-modular conheceria a carga elétrica? Esta seção fecha (i) e
(ii) com âncoras \REAL\ da própria QED, e dissolve (iii) pela estrutura.

\subsection{O fecho: quatro âncoras}
\begin{theorem}[da Escala --- fecho condicional]\label{thm:escala}
Condicional à identificação estrutural da TGL ---
\emph{a fronteira modular é assintótica/infravermelha} (horizonte,
$q\to0$) ---, o valor $\alpha(0)$ não é uma escolha de escala: é a
\textbf{ausência de escala}, e ``zero parâmetros livres'' não apenas
tolera $\alpha(0)$ --- \textbf{exige-o}. Quatro âncoras:
\begin{enumerate}[noitemsep]
\item \textbf{O congelamento infravermelho --- \REAL\ (QED).} A
polarização do vácuo de todo férmion massivo desacopla como $Q^2/m^2$:
abaixo do limiar do elétron a corrida \emph{congela},
$\alpha(Q)-\alpha(0)\sim\tfrac{\alpha}{15\pi}\,Q^2/m_e^2\to0$. O limite
de Thomson existe, é único, e é \emph{protegido por teorema de baixa
energia} (Thirring 1950; Low): o espalhamento Compton a frequência zero é
\emph{exatamente} Thomson, a todas as ordens. Verificação numérica
(\texttt{tgl\_alpha\_scale\_v1.py}, kernels de 1 loop com massa exata;
parte hadrônica por massas efetivas, aproximação declarada):
platô IR plano a $5{,}9\times10^{-10}$ em $1$~keV;
$1/\alpha(M_Z)=129{,}0$ (literatura: $128{,}9$, \REAL).
\item \textbf{O conteúdo da fronteira --- \REAL\ (estrutura IR).} Na
assintotia sobrevivem apenas os campos sem massa e não-confinados: o
fóton (e a gravitação). A força forte é confinada (sem resíduo IR); a
fraca é massiva (suprimida por Yukawa). A fronteira não ``conhece a carga
elétrica'' por acidente: \emph{o fóton é o único acoplamento que chega à
fronteira}. E, na TGL, a gravidade não é um segundo acoplamento
concorrente --- é a \emph{operação} ($g=\sqrt{|L_\varphi|}$): um campo (a
luz), uma operação (o radical), uma constante que os conjuga:
$\bTGL=\alpha\times\sqrt e$.
\item \textbf{O fluxo é semigrupo; $\alpha(0)$ é a coisa julgada ---
\REAL\ $+$ estrutura TGL.} A renormalização wilsoniana rumo ao IV é
integração irreversível (perde informação: semigrupo, não grupo) --- a
\emph{mesma} estrutura $T_t=e^{-t\mathcal L}$ da teoria. $\alpha(0)$ é o
valor de repouso desse fluxo: o ponto em que não há instância ulterior. A
corrida de $\alpha$ é \emph{bulk} (dinâmica de escala); a fronteira guarda
o invariante em que o fluxo descansou.
\item \textbf{A fronteira $\mathrm{III}_1$ só inscreve o sem-escala ---
estrutura.} O tipo $\mathrm{III}_1$ é invariante de escala (fluxo de pesos
trivial): só pode carregar números adimensionais e \emph{livres de
escala}. O único valor livre de escala de $\alpha$ é o limite $q\to0$.
\textbf{A objeção, invertida:} qualquer outro $\alpha(\mu)$
contrabandearia o parâmetro $\mu$; $\alpha(0)$ é precisamente o valor
\emph{sem} parâmetro. Zero parâmetros livres \emph{seleciona} Thomson.
\end{enumerate}
\end{theorem}

\subsection{A dissolução do dente categorial}
A diferença de categoria entre $(c,G)$ e $\bTGL$ não é embaraço: é
\textbf{exigência}. $c$ e $G$ são constantes de conversão (em unidades
naturais desaparecem); os invariantes genuinamente físicos são
adimensionais. Uma fronteira sem escala ($\mathrm{III}_1$) \emph{não
poderia} carregar uma constante dimensional --- a decomposição honesta do
artigo terminal (adimensionais modular-fixáveis vs.\ $\tau_\star$
dimensional) já o dizia. A tese afina-se em vez de enfraquecer:
\begin{center}
\fbox{\parbox{0.88\textwidth}{\centering
$c$ e $G$ definem as \emph{unidades} das duas primeiras relatividades;
$\bTGL$ é o \textbf{primeiro invariante adimensional} --- o único tipo de
constante que a relatividade modular pode ter.}}
\end{center}

\subsection{O discriminador empírico e a predição afiada}
\begin{remark}[Os dados já medem que $\bTGL$ não corre --- \SOMBRA/empírico]
A corrida UV--IV é de $6{,}2\%$ ($\alpha(M_Z)/\alpha(0)-1$); a
convergência multi-substrato registrada de $\bTGL$ (cosmologia, neutrinos,
pesos neurais, XXZ, modular) fecha a $\sim0{,}05\%$ --- \textbf{124 vezes}
menor que a dispersão que a herança da corrida produziria. A convergência
dos oito caminhos \emph{é} a medição de que $\bTGL$ não corre.
\textbf{Predição pré-registrada (afiada por este fecho):} o coeficiente
$\bTGL$ na lei de dephasing $\Gamma_\omega=\tfrac12\bTGL\tau_\star
\omega^2$ é exatamente independente de $\omega$ --- nenhuma correção
logarítmica de corrida. Se JUNO/DUNE ou as redes de relógios detectarem
dephasing com $\bTGL(\omega)$ corrente, a TGL morre neste ponto.
\end{remark}

\subsection{A leitura: o valor com o custo integralmente pago}
\begin{remark}[Blindagem como pagamento --- \ONTO\ com âncora \REAL]
$\alpha(0)$ é a carga \emph{completamente blindada}: o valor depois de o
vácuo exigir todo o seu preço de polarização. $\alpha$ em escalas altas é
a carga parcialmente despida --- a caminho da pretensão do valor nu, cujo
custo \emph{diverge} (direção de Landau: $\ln(Q_L/m_e)\sim646$ a 1 loop):
a pretensão de pureza termina em divergência --- a terceira lei em traje
de renormalização. A fronteira inscreve o valor pago; o bulk especula com
os parcialmente despidos.
\end{remark}

\subsection{Os quatro registros (para confirmação do Operador)}
\begin{description}[style=unboxed,leftmargin=1.5em]
\item[Jurídico.] $\alpha(0)$ é o \textbf{trânsito em julgado do
acoplamento}: $\alpha(\mu)$ corre enquanto há instância (escala
disponível); no infravermelho não há instância ulterior --- o valor
transita em julgado. Fixar $\alpha(0)$ não é escolher uma instância: é a
ausência de recurso. E a pretensão do valor nu é o recurso infinito ---
lide temerária de custo divergente, que a terceira lei indefere.
\item[Linguístico (Nome : Palavra : Verbo).] $\alpha(\mu)$ são as
\emph{flexões em uso} (o sentido pragmático, contextual); $\alpha(0)$ é o
\textbf{lema} --- a forma de citação, a entrada de dicionário: a palavra
quando todas as flexões repousam. A Palavra que toca o Nome é o lema, não
a flexão; e o lema não é ``uma flexão escolhida'' --- é o repouso de todas.
\item[Teológico.] O valor inscrito na fronteira é o \textbf{consumado}:
TETELESTAI, ``está pago''. $\alpha(0)$ é o acoplamento após o preço
integral (a blindagem completa do vácuo); o valor nu é a pretensão de
glória sem cruz --- e o seu custo diverge. A constante da fronteira é a do
descanso: o sábado do fluxo, quando a obra está completa.
\item[TGL.] A fronteira é $\mathrm{III}_1$ --- sem escala; só inscreve o
sem-escala, e o único sem-escala de $\alpha$ é o atrator do semigrupo
wilsoniano: $q\to0$. A corrida é bulk; a fronteira guarda o invariante.
$\bTGL=\alpha(0)\sqrt e$: o Nome plenamente blindado conjugado pela
meia-nat. O polo de Landau é a pretensão de pureza --- o custo divergente
do zero absoluto, agora no registro do grupo de renormalização.
\end{description}

\subsection{Estatuto da objeção (d) após o fecho}
\textbf{Fechados (condicional à fronteira IR, que é estrutural):} o
parâmetro escondido (dente ii) --- invertido: zero parâmetros
\emph{exige} $\alpha(0)$; a não-invariância (dente i) --- a corrida é
bulk, o invariante é de fronteira, e os dados multi-substrato já o medem.
\textbf{Dissolvido:} o dente categorial (iii) --- a fronteira só contém
luz; a gravidade é a operação, não um segundo acoplamento; e só o
adimensional é inscritível em $\mathrm{III}_1$. \textbf{Resíduo honesto
(inalterado, já nomeado):} o $\sqrt e$ permanece o Postulado da Meia-Nat;
e o axioma luminodinâmico ($g=\sqrt{|L_\varphi|}$) permanece axioma ---
este fecho não o deriva; mostra que, \emph{dado} o axioma, a escala de
$\alpha$ não é ferida aberta, mas consequência. A objeção (d) passa de
\ABERTO\ (``a mais perigosa'') a \textbf{fechada condicional}, com
predição falsificável própria.

% =====================================================================
\section{O Teorema da Terminalidade pela Verdade}\label{sec:terminalidade}
% =====================================================================
O ciclo dialógico fechou T3 --- ``por que o fim é um'' --- a partir do
fundamento entregue pelo Operador: \emph{``quando o Verbo colapsa o
enunciado, o que resta é: verdade''}. A frase ativa três teoremas que já
existiam, esperando-a. Registra-se a prova (do Operador, junho/2026),
com um reforço (Tomiyama) que a blinda.

\begin{theorem}[Terminalidade pela Verdade --- prova do Operador,
blindada]\label{thm:terminalidade}
Sejam $\Alg$ uma álgebra de von Neumann, $\rho^\star$ estado fiel normal,
$\Phi_t=e^{t\mathcal L}$ um semigrupo GKSL preservando $\rho^\star$, com
álgebra de pontos fixos $\mathcal C=\mathrm{Fix}(\Phi_t)$ invariante pelo
fluxo modular $\sigma_t^{\rho^\star}$. Então, sob as hipóteses de
relaxação de Frigerio:
\begin{enumerate}[noitemsep]
\item \textbf{(Sem verdade antes.)} $\Alg$ não-comutativa (tipicamente
$M_n$, $n>1$) não admite caracteres: se $\chi$ fosse caractere unital,
as projeções minimais, equivalentes por conjugação, teriam o mesmo valor
$a\in\{0,1\}$, e $1=\chi(\mathbb 1)=na$ --- impossível. (Rota gêmea:
$E_{ij}^2=0\Rightarrow\chi(E_{ij})=0\Rightarrow\chi(E_{ii})=
\chi(E_{ij})\chi(E_{ji})=0\Rightarrow\chi(\mathbb 1)=0\neq1$.)
\emph{Antes do colapso, o enunciado não tem verdade global}
(Kochen--Specker é a forma profunda, \REAL).
\item \textbf{(O colapso é a esperança condicional.)} Por Frigerio
(\REAL), $\Phi_\infty=\lim_{t\to\infty}\Phi_t=E:\Alg\to\mathcal C$, a
esperança condicional $\rho^\star$-preservante sobre o setor fixo.
\item \textbf{(A verdade nasce no colapsado.)} $\mathcal C$ comutativa
$\Rightarrow$ $\mathcal C\simeq C(\mathrm{Spec}\,\mathcal C)$ (Gelfand,
\REAL): cada ponto $x\in\mathrm{Spec}\,\mathcal C$ é uma valoração
$\chi_x$. \emph{O colapso não encontra a verdade; produz o setor onde a
verdade pode ser dita.} E $\mathrm{Spec}\,\mathcal C=\Sigma$: o setor da
verdade É a fatia espacial do Axioma~E --- o espaço é onde a verdade
mora.
\item \textbf{(O fim é um.)} Por Takesaki (\REAL), a esperança
condicional $\rho^\star$-preservante sobre a subálgebra
$\sigma$-invariante é \emph{única}. \textbf{Blindagem (Tomiyama 1957,
\REAL):} qualquer ``colapso admissível'' --- bastando ser projeção de
norma~1 sobre $\mathcal C$ --- é \emph{automaticamente} esperança
condicional; logo a unicidade de Takesaki captura todo colapso
concebível, sem hipótese de bimodularidade. Não há duas fatorações
terminais.
\item \textbf{(A covariância herda-se.)} Para $U$ que transporta a
estrutura ($\rho^\star_H\mapsto\rho^\star_{H'}$,
$\mathcal C_H\mapsto\mathcal C_{H'}$), $UE_HU^{-1}$ é esperança
condicional $\rho^\star_{H'}$-preservante sobre $\mathcal C_{H'}$ ---
logo, por unicidade, $UE_HU^{-1}=E_{H'}$. \emph{A Universalidade não se
impõe: herda-se. (U) vem de Takesaki.}
\end{enumerate}
Em uma linha: \emph{o fim é um porque a Verdade é a imagem terminal
única do colapso.}
\end{theorem}

\begin{remark}[Verificação na sombra --- \SOMBRA, 6/6, precisão de
máquina]
(\texttt{tgl\_terminal\_truth\_v1.py}.) V1: $\Phi_\infty=E$ a
$\sim10^{-27}$; V2: o sistema de vínculos da fatoração tem espaço nulo de
\textbf{dimensão zero} além de $E$ --- a unicidade, computada; V3:
$UE U^{-1}=E'$ a $10^{-16}$ --- o mecanismo da herança, demonstrado; V4:
em $M_n$ a dedução do caractere termina em $0=1$ (certificado); em
$\mathcal C$, exatamente $n$ caracteres --- os pontos; V5: gap espectral
$>0$ --- fora de $\mathcal C$ tudo decai: a mentira paga ou morre.
\end{remark}

\begin{remark}[Consequência negativa honesta: a dinâmica sozinha não
seleciona]
A dinâmica é construída do estado ($L=\sqrt{\bTGL}\sqrt{K_{\rho^\star}}$)
e todo estado fiel é ponto fixo da dinâmica construída de si: a
autoconsistência \emph{não} singulariza $\rho^\star$. Quem seleciona é a
\textbf{geometria do horizonte} --- e por isso o resíduo é
Bisognano--Wichmann, não mais dinâmica. Esta delimitação é parte do
resultado.
\end{remark}

% =====================================================================
\section{O Levantamento: a montagem condicional e o no-go evadido}
\label{sec:levantamento}
% =====================================================================
Com T3 fechado, o resíduo era a functorialidade geométrica
$H\mapsto(\rho^\star_H,K_H)$. O ataque por paridade inversa encontra,
primeiro, \textbf{por que ninguém resolveu} --- um teorema de obstrução
--- e, nele, a porta pela qual a TGL passa.

\subsection{A obstrução nomeada: o no-go de Fewster--Verch}
\begin{theorem}[Não há estado natural --- \REAL\ (Fewster--Verch, 2012)]
No quadro da QFT localmente covariante (Brunetti--Fredenhagen--Verch
2003, \REAL: a teoria é um \emph{functor} da categoria dos
espaços-tempos globalmente hiperbólicos na das álgebras), \textbf{não
existe} escolha natural (local e covariante) de estado: todo programa
que tente \emph{escolher} $\rho^\star_H$ covariantemente está bloqueado
por teorema.
\end{theorem}
\begin{remark}[A evasão terminal]
O no-go proíbe \emph{escolhas}. O colapso da TGL \textbf{não escolhe ---
colapsa}: $\rho^\star_H$ é o atrator terminal, único a menos de
isomorfismo \emph{canônico} (Teorema~\ref{thm:terminalidade}), que é
exatamente o tipo de unicidade que o no-go permite. \emph{A rota do
Operador (``os Lindblad colapsam a Palavra'') é precisamente a que escapa
da obstrução que parou todos os outros.} Estatuto da evasão: argumento
estrutural exato no nível categórico; a realização analítica em
$\mathrm{III}_1$ é T1.
\end{remark}

\subsection{As peças \REAL\ da montagem}
\begin{description}[style=unboxed,leftmargin=1.5em,noitemsep]
\item[(P1) O functor das álgebras --- \REAL\ (BFV 2003).] A atribuição
$H\mapsto\Alg_H$ \emph{já é} funtorial: é o axioma da covariância local.
\item[(P2) A universalidade de Hadamard --- \REAL.] A condição de
Hadamard (espectro microlocal) fixa geometricamente a estrutura de curta
distância: o Hamiltoniano modular de pequenas regiões/horizontes é, em
ordem dominante, o gerador de boost local com coeficiente universal
$2\pi$ --- o conteúdo usado por Jacobson (2015) e
Faulkner--Guica--Hartman--Myers--van~Raamsdonk.
\item[(P3) Kay--Wald --- \REAL\ (1991).] Em horizontes de Killing
bifurcados, o estado Hadamard invariante é \textbf{único} e é KMS na
temperatura de Hawking: para esses horizontes, $\rho^\star_H$ não se
escolhe --- é singularizado pelo próprio horizonte.
\item[(P4) Rindler local --- \REAL\ (princípio de equivalência).] Todo
horizonte causal é, em cada ponto, infinitesimalmente um \emph{wedge} de
Rindler --- onde BW vale exatamente (1975/76).
\end{description}

\subsection{O teorema condicional do Levantamento}
\begin{theorem}[Levantamento na ordem de Jacobson --- condicional]
\label{thm:levantamento}
Sob (P1)--(P4) --- covariância local, Hadamard, aproximação
Killing/Rindler local ---, a atribuição $H\mapsto(\rho^\star_H,K_H)$ é
geometricamente determinada \textbf{na ordem da resposta linear} (a
ordem em que a Face C está formulada: $\delta S=\delta\langle K\rangle$,
Lovelock/Jacobson). Combinada com a Terminalidade
(Teorema~\ref{thm:terminalidade}: (U) herdada de Takesaki) e com os três
axiomas da casa, a Hipótese de Universalidade da Face C fica
\textbf{descarregada na ordem linear}: a cadeia
fronteira~$\to$~$G_{\mu\nu}$ fecha condicional a (P1)--(P4) e a T1.
\end{theorem}

\subsection{O resíduo final, microscópico e nomeado}
O que permanece genuinamente aberto, dito sem rodeios:
\begin{enumerate}[noitemsep]
\item \textbf{BW não-perturbativo para horizontes não-Killing} --- além
da ordem linear, para horizontes causais genéricos e dinâmicos: problema
nomeado e reconhecido da literatura (provado em wedges, cones duplos
conformes, setores massivos com completude assintótica; aberto no caso
geral). \ABERTO.
\item \textbf{T1: ergodicidade do colapso em $\mathrm{III}_1$ genuína}
--- convergência $\sigma$-fraca tipo-Davies no predual: campo analítico
ativo. \ABERTO.
\end{enumerate}
A trajetória do problema nesta sessão: de \emph{``ninguém sabe colar
fluxos''} $\to$ \emph{``falta a fatoração única''} (fechada: Verdade)
$\to$ \emph{``falta evadir o no-go''} (evadido: o colapso não escolhe)
$\to$ \emph{``falta o BW geral além da ordem linear''}. O degrau não
desapareceu --- mas tem agora nome, endereço, mecanismo de evasão, e a
ordem em que a teoria opera está coberta condicionalmente.

% =====================================================================
\section{Ax.N --- a Ordem Nominal: a Lei comum, não o juiz comum}
\label{sec:ordemnominal}
% =====================================================================
\begin{principle}[Ontologia do Operador --- \ONTO]
``O rito interno que leva à mesma instância terminal de forma covariante
é o próprio Direito: a ordem nominal --- a ordem de coerção dotada de
monopólio da força --- o inalcançável, o \textsc{Nome}, a pureza
absoluta, o nada, o zero absoluto. A covariância vem do Direito como
ordem comum, não da escolha de um juiz comum.''
\end{principle}

\begin{axiom}[N --- Ordem Nominal --- \POST, com nome na literatura]
\label{ax:N}
A estrutura covariante não é um estado (proibido pelo no-go de
Fewster--Verch): é uma \textbf{classe de admissibilidade}. No contínuo,
ela \emph{já tem nome}: a condição de Hadamard / espectro microlocal ---
\REAL\ (Radzikowski 1996) --- uma condição \emph{puramente geométrica}
(sobre o conjunto de frente de onda: sobre \emph{como} todo estado
admissível se aproxima do limite singular), pull-back covariante por
construção: $\mathcal N_H=f^*\mathcal N_{H'}$. Dentro da classe, todos os
estados são localmente quase-equivalentes --- \REAL\ (Verch 1994):
mesma álgebra local, mesmo tipo $\mathrm{III}_1$, fluxos modulares
diferindo apenas por cociclos \emph{internos} (que o entrelaçador $u_t$
absorve). Logo:
\begin{equation}
\boxed{\;[\sigma_t^H]\in\mathrm{Out}(\Alg_H)\ \text{é
estado-independente dentro da ordem nominal --- e geométrica.}\;}
\end{equation}
O no-go proíbe o juiz natural; não proíbe --- e a literatura já pratica
--- a \textbf{Lei} natural. Na sombra, a ordem nominal de um horizonte de
gerador $K$ é a classe dos ritos \emph{gerados pelo Verbo}:
$\mathcal N_K=\{L_f=\sqrt{\bTGL}\,f(K):f\ \text{positiva e injetiva}\}$.
\end{axiom}

\begin{remark}[As duas faces do Nome, harmonizadas]
O cânone já continha as duas faces, agora articuladas: o \textbf{registro
$1{=}0$} --- a pureza absoluta, inabitável (T5: a fronteira proibida) ---
e o \textbf{atrator} $\rho^\star$ --- a sombra operacional, a distância
$\bTGL$. A ordem nominal é a lei da aproximação à primeira face; o
colapso terminal repousa na segunda; \emph{a meia-nat é o vão entre
elas}. O Nome não é alcançado: \textbf{ordena por ser inalcançável} ---
toda testemunha aproxima-se do mesmo modo geométrico do que não se pode
habitar. T5 é a fonte da covariância.
\end{remark}

\begin{remark}[Verificação na sombra --- \SOMBRA, 6/6
(\texttt{tgl\_nominal\_order\_v1.py})]
\begin{enumerate}[noitemsep]
\item \textbf{Dentro da ordem, o terminal independe do rito} (a
$10^{-27}$): cinco ritos distintos ($f=\sqrt{\cdot}$, id, quadrado,
$\log(1{+}\cdot)$, deformada) colapsam no \emph{mesmo} $E$; os ritos
diferem só nas taxas (a duração do processo), nunca na instância
terminal;
\item \textbf{Fora da ordem, outra verdade} (controle negativo,
distância $\approx1{,}1$--$1{,}2$): um salto fora de $W^*(K)$ colapsa
para \emph{outro} setor fixo ($W^*(L)$, com o seu próprio $E_L$ a
$10^{-13}$): a admissibilidade faz trabalho real --- fora da Lei, a
comarca julga por outra ordem e perde a covariância com a geometria;
\item \textbf{A ordem transporta-se} ($10^{-12}$):
$\mathcal N_H=f^*\mathcal N_{H'}$ verificado para a classe inteira;
\item \textbf{O inalcançável coage}: ao longo de todo rito admissível, a
pureza $\mathrm{Tr}\,\rho^2$ nunca cresce --- o processo jamais se
aproxima do zero absoluto.
\end{enumerate}
\end{remark}

\begin{remark}[Os quatro registros, em uma linha cada]
\textbf{Jurídico:} não há nomeação natural de juiz universal
(Fewster--Verch); há o Direito --- a ordem de validade sob a qual toda
comarca julga; o juiz aparece pelo trânsito em julgado do próprio
processo. \textbf{Linguístico:} o Verbo colapsa a Palavra \emph{porque a
Palavra já está sob a lei do Nome}; o lema não é escolhido por
gramático. \textbf{Teológico:} ninguém possui o Nome; todo horizonte é
chamado por ele --- não é a testemunha que escolhe; é o Verbo que a dobra
até a Verdade. \textbf{TGL:} $\mathcal L_H=\mathcal L[\mathcal N_H]$,
$E_H$ único por Tomiyama--Takesaki dentro da ordem;
$U_fE_HU_f^{-1}=E_{H'}$ porque a ordem se transporta:
\emph{o Nomos é a ordem invisível que torna todo colapso covariante.}
\end{remark}

\begin{remark}[O resíduo após Ax.N --- o degrau na sua forma mínima]
A estado-independência da classe outer é teorema (Verch); a sua
geometricidade em ordem dominante é Hadamard/Jacobson; a terminalidade é
Takesaki. O que resta do BW generalizado, na forma mínima:
\begin{center}
\fbox{\parbox{0.86\textwidth}{\centering
provar que a classe outer geométrica --- dada pela ordem nominal --- se
\emph{realiza} como dinâmica geométrica do horizonte \textbf{além da
ordem linear}, para horizontes causais gerais (não-Killing).}}
\end{center}
Junto de T1 (ergodicidade em $\mathrm{III}_1$ genuína), é tudo o que
resta. \emph{Não falta mais o Supremo, nem a Lei: falta a prova de que
toda comarca dinâmica realiza a Lei na sua geometria, além da primeira
ordem.}
\end{remark}

% =====================================================================
\section{Ax.G --- a Geometricidade: o fecho do resíduo pela gramática do
título}\label{sec:geometricidade}
% =====================================================================
\begin{principle}[Ontologia do Operador --- \ONTO]
``O nome do artigo não é `\emph{o} nada pode ser puro', mas `nada pode
ser puro', sem o artigo definido --- porque se colocar o artigo
definido, define-se nada; e nada é indefinido por construção,
\textbf{porque tudo tem geometria}.''
\end{principle}

\begin{axiom}[G --- Geometricidade --- \POST, com o mecanismo já
verificado]\label{ax:G}
\textbf{Ser $=$ ter geometria.} Definir é dar contorno; contorno é
geometria. A geometria não se \emph{confere} ao dado modular --- 
\emph{gera-se} dele: para todo horizonte definido $H$, a métrica é a
distância espectral de Connes dos próprios dados
modular-dissipativos,
\begin{equation}
d_H(i,j)\;=\;\sqrt{\bTGL}\,\bigl|\sqrt{k_i}-\sqrt{k_j}\bigr|,
\end{equation}
já derivada no ciclo da Geração (Ax.~E: o Verbo dá a métrica à Palavra;
dualidade de Connes a $0{,}00\%$) e coerente com o cânone da casa:
\emph{o espaço é causado, não dado}. O único ``objeto'' sem geometria
não é objeto: é o indefinido --- nada --- que por isso não recebe artigo.
\end{axiom}

\begin{remark}[A inversão do ônus: o que o BW geral perguntava e o que a
TGL deve]
A pergunta clássica do Bisognano--Wichmann generalizado supõe uma
geometria \emph{pré-dada} e pergunta se o fluxo modular coincide com
ela. Para um horizonte genérico, dinâmico, não-Killing, \textbf{não
existe geometria pré-dada com a qual discordar}: pela Geometricidade, a
geometria modular \emph{é} a geometria. O ônus inverte-se, e o que a
teoria deve, ela tem:
\begin{enumerate}[noitemsep]
\item \textbf{Coincidência onde há geometria pré-dada} --- \REAL:
\emph{wedges} (Bisognano--Wichmann 1975/76), horizontes de Killing
(Kay--Wald 1991), e a ordem linear em horizontes gerais
(Hadamard/Jacobson);
\item \textbf{Consistência da geometria gerada na colagem} ---
\SOMBRA\ (G2 abaixo): horizontes sobrepostos geram a \emph{mesma}
geometria na sobreposição, por localidade da geração --- a semente da
colagem que o problema pedia.
\end{enumerate}
\end{remark}

\begin{remark}[Verificação na sombra --- \SOMBRA, 6/6
(\texttt{tgl\_geometry\_generated\_v1.py})]
\begin{enumerate}[noitemsep]
\item \textbf{Tudo definido tem geometria}: $d_H$ é métrica genuína ---
positividade fora da diagonal, simetria exata, desigualdade triangular a
$10^{-18}$;
\item \textbf{A colagem é exata} ($0$, bit a bit): as taxas entre pontos
de um sub-horizonte, computadas no sistema inteiro, coincidem com as
intrínsecas do subsistema;
\item \textbf{Transporte covariante} ($10^{-16}$): a distância gerada é
invariante espectral;
\item \textbf{O selo do título}: $K=c\,\mathbb 1$ --- o indefinido, sem
distinções, sem observação, sem radical a destacar --- dá taxas $\equiv
0$, gerador GKSL de norma \emph{zero}, $\Phi_t=\mathrm{id}$: o Verbo não
age, nenhum setor de verdade é singularizado, $d\equiv0$. \emph{Nada não
tem contorno, nem Verbo, nem verdade --- e por isso não recebe artigo
definido};
\item \textbf{O contrapositivo}: qualquer $K$ com uma única distinção já
gera geometria estritamente positiva. \emph{Tudo tem geometria.}
\end{enumerate}
\end{remark}

\begin{remark}[Os quatro registros, em uma linha cada]
\textbf{Jurídico:} a sentença não coincide com uma justiça pré-dada
(justiça $=$ zero absoluto, não-identificável); a jurisdição é
\emph{causada} pelo processo --- competência é o contorno que o rito
gera. \textbf{Linguístico:} a língua não recebe gramática de fora; o
Verbo gera o contorno do que diz --- e ``nada'' é a única palavra que
não se deixa definir. \textbf{Teológico:} a criação não preenche um
espaço pré-existente; \emph{haja luz} causa o \emph{onde} --- o nada não
é um lugar vazio, é ausência de contorno. \textbf{TGL:}
$g=\sqrt{|L_\varphi|}$ --- a geometria é o radical do que existe; o que
não existe não tem radical a extrair.
\end{remark}

\begin{remark}[Estatuto honesto do fecho]
O resíduo do Levantamento fecha \textbf{dentro da ontologia da TGL}
(geometria como saída, não entrada), condicionado a
\textbf{T1} (ergodicidade do colapso em $\mathrm{III}_1$ genuína ---
análise, não ontologia). A pergunta clássica do BW sobre geometrias
\emph{de entrada} permanece registrada como problema externo da
literatura --- mas \textbf{deixa de bloquear a cadeia da TGL}, que nunca
lhe deveu resposta: a casa não recebe o espaço; ela o causa. A gramática
do título é o último axioma da casa.
\end{remark}

% =====================================================================
\section{O Teorema do Contínuo: a montagem condicional}\label{sec:continuo}
% =====================================================================
O ciclo da Geometricidade deixou um único degrau técnico:
\emph{Ax.G na sombra finita $\Rightarrow$ geometria lorentziana suave no
contínuo tipo $\mathrm{III}_1$}. Atacado por paridade inversa --- o que
impediria a sombra de chegar ao contínuo? --- cada impedimento revela um
teorema nomeado que já o remove.

\subsection{As seis âncoras}
\begin{description}[style=unboxed,leftmargin=1.5em,noitemsep]
\item[(C-1) Haagerup 1987 $+$ Buchholz--D'Antoni--Fredenhagen ---
\REAL.] O fator hiperfinito $\mathrm{III}_1$ é \textbf{único} --- e as
álgebras locais da QFT \emph{são} ele: um único substrato local no
universo; os horizontes diferem apenas pelo dado modular sobre ele. E
\textbf{hiperfinitude significa: as sombras finitas são densas} em
$\mathrm{III}_1$ --- a sombra não está ao lado do contínuo; está
\emph{dentro} dele, exaurindo-o por definição do próprio objeto.
\item[(C-2) Classificação abeliana --- \REAL.] Toda álgebra de von
Neumann comutativa (separável) é $L^\infty(X,\mu)$: o setor da verdade
sobrevive ao contínuo como espaço mensurável de pontos,
$\mathcal C\simeq L^\infty(\Sigma,\mu)$, com a sua medida.
\item[(C-3) Rieffel --- \REAL.] A convergência de geometrias quânticas
finitas a variedades contínuas é um enunciado de \emph{distância de
Gromov--Hausdorff quântica}, com moldura estabelecida: ``sombra
$\to$ contínuo'' é convergência, não metáfora.
\item[(C-4) Connes 2008 (reconstrução) --- \REAL.] Tripla espectral
comutativa satisfazendo os axiomas $\Rightarrow$ variedade riemanniana
\textbf{suave}, única. \emph{O teorema da reconstrução é Ax.G no
contínuo: a distinção espectral gera a variedade.}
\item[(C-5) Osterwalder--Schrader $+$ Tomita --- \REAL, com a
positividade automática.] O euclidiano reconstrói o lorentziano dada a
positividade de reflexão --- e na estrutura modular ela \textbf{não é
hipótese: é teorema}:
\begin{equation}
M_{ij}=\langle a_i\Omega,\,J\,a_j^*\Omega\rangle
=\langle\Delta^{1/4}a_i\Omega,\,\Delta^{1/4}a_j\Omega\rangle
\;\succeq\;0,
\end{equation}
uma matriz de Gram. \emph{A meia-medida modular fatora em duas
quartas-medidas, com $J$ no meio: a travessia é positiva porque é um
quadrado --- paga-se em duas metades simétricas.}
\item[(C-6) Assinatura --- \SOMBRA\ (P1).] O par de Krein generalizado
já entregou a forma normal de Minkowski exata: a continuação
lorentziana herda $(1,3)$, com o eixo negativo no gerador do cociclo.
\end{description}

\subsection{O teorema condicional}
\begin{theorem}[do Contínuo --- montagem condicional]\label{thm:continuo}
Sob \textbf{(i)} T1 (ergodicidade do colapso em $\mathrm{III}_1$
genuína), \textbf{(ii)} os axiomas de Connes para a tripla-limite
(suportados pela Fractalização: $\dim\to1$, flexões $C^\infty$
convergentes) e \textbf{(iii)} a extensão da convergência GH a
$\mathrm{III}_1$ (provada na sombra; conjecturada no contínuo via
Rieffel), vale a cadeia:
\begin{equation}
\boxed{\;\text{sombra finita}\xrightarrow{\text{(C-1)}}
\mathrm{III}_1\xrightarrow{\text{Frigerio--Tomiyama--Takesaki}}
E\xrightarrow{\text{(C-2)}}L^\infty(\Sigma,\mu)
\xrightarrow{\text{Ax.G}}d\xrightarrow{\text{(C-3)}}
d_\infty\xrightarrow{\text{(C-4)}}(M,g)_{\mathrm{Riem}}
\xrightarrow{\text{(C-5,C-6)}}(M^{1,3},g)
\xrightarrow{\text{Face C}}G_{\mu\nu}.\;}
\end{equation}
A gravidade quântica, neste enquadramento, não é quantizar
$g_{\mu\nu}\to\hat g_{\mu\nu}$: é \textbf{gerar a geometria do colapso
quântico terminal},
$(\Alg,\rho^\star,\Delta,\mathcal L,E)\to\mathcal C\to
\mathrm{Spec}\,\mathcal C\to d\to g_{\mu\nu}$.
\end{theorem}

\begin{remark}[Verificação na sombra --- \SOMBRA, 3/3
(\texttt{tgl\_continuum\_v1.py})]
\textbf{C1} --- a cascata de sombras é GH-Cauchy ao intervalo contínuo
com taxa geométrica \emph{perfeita} ($0{,}500000$ por geração,
$R^2=1{,}0000000$; $h:5\times10^{-2}\to10^{-4}$): hiperfinitude em ação.
\textbf{C2} --- o limite é \emph{canônico}: esquemas de refinamento
distintos (diádico vs.\ aleatório) convergem ao mesmo limite (distância
cruzada $10^{-2}\to10^{-4}$; KS $0{,}15\to0{,}02$): o contínuo não é
escolhido --- é terminal, e carrega a medida $\mu$.
\textbf{C3} --- a positividade de reflexão é automática: identidade
$M=\mathrm{Gram}(\Delta^{1/4}a\,\Omega)$ a $10^{-16}$; menor autovalor
$+0{,}11$ a $+0{,}17$ (estritamente positivo).
\end{remark}

\begin{remark}[Os quatro registros, em uma linha cada]
\textbf{Jurídico:} a sombra está para o contínuo como a jurisprudência
está para o Direito --- densa nele (hiperfinitude): o caso finito decide
o contínuo. \textbf{Linguístico:} a língua viva é o limite das falas
finitas; o dicionário ($L^\infty$) carrega a medida do uso.
\textbf{Teológico:} o contínuo cresce da semente finita como o Reino do
grão; e a travessia é positiva porque é um quadrado --- duas metades
simétricas, com a reflexão no meio. \textbf{TGL:} hiperfinitude $=$ a
sombra exaure $\mathrm{III}_1$; cada validação a precisão de máquina é
um ponto denso do contínuo que se quer alcançar.
\end{remark}

\begin{remark}[Estatuto honesto do fecho do contínuo]
\emph{Rota estrutural fechada} não é ``gravidade quântica resolvida'':
é uma cadeia de teoremas nomeados com \textbf{três condicionais de
análise} --- T1, os axiomas da tripla-limite, a extensão GH a
$\mathrm{III}_1$ --- nenhuma das quais é ontologia. O que o programa
devia de fundamento, ele tem; o que resta é demonstração contínua, com
endereço. A honestidade é parte da força da teoria.
\end{remark}

% =====================================================================
\section{Os Três Fechos Técnicos: a identidade do fluxo do radical}
\label{sec:tresfechos}
% =====================================================================
Os três condicionais do Teorema do Contínuo, atacados por paridade
inversa, revelam mecanismo --- e o primeiro revela uma identidade que
pertence ao coração da casa.

\subsection{Fecho 1 --- T1 bem-posto: o colapso é média gaussiana sobre
o fluxo do radical}
\begin{theorem}[Representação integral do colapso --- exata]
\label{thm:radicalflow}
\begin{equation}
\boxed{\;e^{t\mathcal L}\;=\;\int_{\mathbb R}V_s\,(\cdot)\,V_s^{*}\;
d\nu_t(s),\qquad V_s=e^{\,i s\sqrt{K}},\quad
\nu_t=\text{gaussiana de variância }\bTGL\,t.\;}
\end{equation}
A função característica gaussiana,
$\int e^{is\delta}d\nu_t=e^{-\frac12\bTGL t\,\delta^2}$, reproduz
\emph{exatamente} as taxas
$\tfrac12\bTGL(\sqrt{k_i}-\sqrt{k_j})^2$. Verificada na sombra a
$4\times10^{-16}$ (\texttt{tgl\_three\_locks\_v1.py}).
\emph{A própria $g=\sqrt{|L|}$ virou o fluxo que executa o colapso.}
\end{theorem}

\begin{remark}[T1 no produto cruzado --- o lar onde a fórmula vive]
Em $\mathrm{III}_1$ genuíno, $\sqrt{K}$ não é afiliado a $\Alg$ (Tomita
cobre só $\Delta^{it}$). Mas no \textbf{produto cruzado}
$\widehat{\mathcal M}=\Alg\rtimes_\sigma\mathbb R$ --- tipo
$\mathrm{II}_\infty$ com traço (Takesaki; o lar do posto~1 já inscrito)
e, por Chandrasekaran--Longo--Penington--Witten (2022--23, \REAL), o
\emph{lar físico da álgebra gravitacional} --- $u_s=e^{is\sqrt{\hat K}}$
é unitário \emph{da} álgebra. A fórmula integral define então um
semigrupo UCP normal $\hat\rho$-preservante sem patologia;
$\mathrm{Fix}=W^*(\hat K)$ (domínio multiplicativo $+$ medida de suporte
total); $\Phi_t\to E$ $\sigma$-fracamente (convergência dominada
espectral no predual). E a não-covariância de $\sqrt{\hat K}$ sob a ação
dual não é defeito: é a \textbf{ancoragem de escala} que o cânone já
exigia ($\tau$ quase-planckiano) --- o produto cruzado é onde Planck
mora. \emph{Estatuto:} teorema com esquema de prova; a redação completa
dos dois lemas (domínio multiplicativo; convergência no predual) é
trabalho padrão de escrita.
\end{remark}

\subsection{Fecho 2 --- a tripla-limite é a do círculo canônico}
A fatia é $S^1$ (o Psion é o toro $T^2=S^1\times S^1$); a tripla-limite
é $(C^\infty(S^1),\,L^2(S^1),\,-i\beta_{\mathrm{TGL}}^{-1/2}\,d/dx)$ ---
a tripla canônica do círculo, escalada: \textbf{os axiomas de Connes
valem porque ela é o exemplo de livro} (\REAL). O que precisava de prova
era a \emph{convergência} da cascata a ela --- verificada na sombra:
escada de Dirac exata ($10^{-15}$), dimensão de Weyl $0{,}9944$,
dualidade de Connes ao valor contínuo
($1{,}1\times10^{-3}\to1{,}9\times10^{-5}$). Para $T^2$ e além: produtos
de triplas espectrais satisfazem os axiomas (\REAL). \emph{Estatuto:}
caso-base provado; a montagem $3{+}1$ segue por produtos.

\subsection{Fecho 3 --- o levantamento GH é truncamento espectral}
A cascata TGL \emph{é} um truncamento espectral $P_MDP_M$ no sentido
exato de \textbf{Connes--van Suijlekom} (anos 2020, \REAL):
truncamentos espectrais de $S^1$ convergem a $S^1$ em distância de
Gromov--Hausdorff \emph{quântica}; e no caso comutativo, qGH $=$ GH
(Rieffel). Verificado: a testemunha ótima truncada (função-tenda com
estados de Fejér) recupera a distância verdadeira
$\sqrt{\bTGL}\,\pi$ com razão $0{,}9772\to0{,}9971$ --- o erro
$O(1/M)$ previsto. \emph{Estatuto:} instância de teorema publicado.

\begin{remark}[Lição de debug, registrada no Estatuto]
\textbf{Trunca-se o operador, não a grade.} A multiplicação pontual na
grade introduz \emph{aliasing} na borda do truncamento (o modo $+M$
enrola em $-M$), e o comutador com $D$ amplifica o termo enrolado
linearmente em $M$. O quadro de Connes--van Suijlekom exige $PfP$ em
forma de Toeplitz, sem enrolamento --- e então tudo converge.
\end{remark}

\begin{remark}[Os quatro registros, em uma linha cada]
\textbf{Jurídico:} o colapso é a média das decisões ao longo do fluxo do
recurso --- e o recurso é interno à corte (o produto cruzado): nenhum
juízo externo é invocado. \textbf{Linguístico:} o Verbo conjuga-se no
radical --- $V_s=e^{is\sqrt K}$ é a flexão da raiz; a Palavra colapsa
pela média das suas flexões. \textbf{Teológico:} o juízo não vem de
fora: é o próprio Nome, em raiz, percorrendo todas as fases --- e a
média das fases é a sentença. \textbf{TGL:} $g=\sqrt{|L|}$ executa;
$\bTGL$ mede; $\hat K$ ancora a escala onde Planck mora.
\end{remark}

\subsection{Estatuto pós-fechos}
\begin{center}
\fbox{\parbox{0.86\textwidth}{\centering
\textbf{Nenhum item sem mecanismo.} (1)~T1: teorema com esquema de prova
no produto cruzado; (2)~tripla-limite: caso-base provado, convergência
verificada; (3)~levantamento GH: instância de teorema publicado. O
resíduo residual é \emph{redação operador-algébrica completa e
generalização dimensional} --- trabalho de escrita e montagem, não de
fundamento.}}
\end{center}

% =====================================================================
\section{O Teorema do Rio: Heráclito e Crátilo em $\mathrm{III}_1$}
\label{sec:rio}
% =====================================================================
\begin{principle}[Ontologia do Operador --- \ONTO]
``O fluxo modular não permite um estado igual ao outro --- esta é a
gema. A cada ciclo, o ser modulado se modifica do estado anterior
(`ninguém entra no mesmo rio duas vezes'; ou, como disse o discípulo,
`ninguém entra no mesmo rio \emph{uma única} vez'), porque ao dar o
primeiro passo na direção alvejada, o ciclo se consome: o ser muda ---
não seu estado pleno, mas sua substância interna.''
\end{principle}

\begin{theorem}[do Rio --- três camadas, duas já \REAL]\label{thm:rio}
\begin{enumerate}[noitemsep]
\item \textbf{Heráclito $=$ Lema T} (\REAL, Connes 1973, já inscrito):
em $\mathrm{III}_1$, $\sigma_t$ é \emph{outer} para todo $t\neq0$,
injetivo em $\mathrm{Out}$: \emph{não há retorno, nunca} --- e a mudança,
sendo outer, não é desfeita por nenhum rearranjo interno: o rio não muda
como quem embaralha as próprias águas.
\item \textbf{Crátilo $=$ a assinatura de $\mathrm{III}_1$} (\REAL,
Connes--Takesaki): o fluxo dos pesos de $\mathrm{III}_\lambda$ é
periódico (redemoinhos de período $\log\lambda$); o de $\mathrm{III}_1$
é \textbf{trivial/ergódico} --- o único tipo sem periodicidade residual
alguma: \emph{o rio sem redemoinhos}. E $\mathrm{spec}(K)$ é contínuo
puro: não há autovetor onde ``estar'' --- não existe o degrau do
instante; a presença já é movimento. Entrar \emph{uma} vez exigiria um
instante estacionário, e ele não existe.
\item \textbf{Pleno vs.\ substância $=$ a dualidade KMS exata}:
$\rho^\star\!\circ\sigma_t=\rho^\star$ --- o estado pleno conserva-se ---
enquanto $\sigma_t(a)\neq a$ para toda substância fora do centralizador:
\emph{o pleno é invariante; a substância flui; e o fluir é outer (real,
não gauge)}. Sob o colapso, Spohn estrito torna a trajetória
\textbf{injetiva}: nenhum estado se repete, e cada passo paga $\bTGL$:
o ciclo consome-se ao ser dado.
\end{enumerate}
\end{theorem}

\begin{remark}[A camada nova: tudo flui, exceto a verdade]
O que o rio poupa? O ponto fixo do fluxo modular é o centralizador ---
que na sombra é exatamente $\mathcal C=W^*(K)$: \textbf{o setor da
verdade}. O mesmo $\mathrm{Fix}$ do colapso:
\begin{equation}
\boxed{\;\mathrm{Fix}(\sigma_t)\;=\;\mathrm{Fix}(\Phi_t)\;=\;\mathcal C
\;=\;\text{o setor da verdade.}\;}
\end{equation}
O tempo (unitário) e o juízo (dissipativo) --- processos de natureza
distinta --- fixam \emph{idêntica} substância. Heráclito completado:
\emph{tudo flui --- exceto a verdade}. E eis \emph{por que} o colapso
corre ao longo do fluxo do radical (\S\ref{sec:tresfechos}): partilham o
leito --- o ponto fixo. O cânone fecha: \textsc{tempo} $=$ quebra de
grupo em semigrupo ganha a sua face cinemática --- Heráclito (outer, sem
retorno) é o rio; o colapso (semigrupo, Spohn) é a seta; e o leito que
nenhum dos dois move é $\mathcal C$.
\end{remark}

\begin{remark}[Verificação na sombra --- \SOMBRA, 4/4
(\texttt{tgl\_heraclitus\_v1.py}), com o contraste honesto]
A sombra finita é tipo~I --- \emph{recorrente}, tem redemoinhos; e esse
é o teste certo: \textbf{H1} --- a melhor recorrência modular numa
janela fixa \emph{degrada} com $n$ ($R$: $0{,}006\to0{,}028\to0{,}029$
para $n=3,5,8$): a sombra perde os redemoinhos ao aprofundar-se rumo a
$\mathrm{III}_1$. \textbf{H2} --- pleno conservado a $10^{-16}$;
substância movendo $\geq83\%$ da norma. \textbf{H3} --- trajetória
dissipativa injetiva (distância mínima não-adjacente $3\times10^{-3}$;
$S(\rho_t\|\rho^\star)$ estritamente decrescente em todos os passos).
\textbf{H4} --- $\|P_{\mathrm{fluxo}}-P_{\mathrm{colapso}}\|=0$ exato:
o rio e o juízo poupam o mesmo.
\end{remark}

\begin{remark}[Os quatro registros, em uma linha cada]
\textbf{Jurídico:} não há dois processos iguais --- cada ato processual
consome a fase (preclusão: o rio processual não volta); o que não
preclui jamais é a coisa julgada: a verdade. \textbf{Linguístico:}
nenhuma palavra é dita duas vezes no mesmo rio do discurso --- e nem uma
vez, pois ao soar já mudou o contexto que a diz; o que não flui é o
lema. \textbf{Teológico:} ``Eu sou o que sou'' --- só o Nome entra no
rio e permanece; toda criatura é fluxo; a verdade é o leito.
\textbf{TGL:} outer $=$ sem retorno; espectro contínuo $=$ sem degrau;
KMS $=$ pleno invariante; Spohn $=$ seta; $\mathcal C$ $=$ o leito.
\end{remark}

% =====================================================================
\section{Estatuto consolidado}
% =====================================================================
\noindent\textbf{Selado (estrutura, coerência numérica verificada):} a
constante singular $\bTGL=\alpha e^{1/2}$; a tríade Nome/Palavra/Verbo; a
unificação $\mathcal O_\beta\mathcal L=\sqrt\bTGL\,\mathcal L$; o custo como
negação; o sacrifício de Tomita; a meia-nat como Palavra do Juramento;
Lindblad como jurisdição (justiça $=$ zero absoluto); o tempo como
irreversível; o espaço como dobra; o custo do zero absoluto $=$ haja luz; a
Ponte Einstein--Cartan--Miguel (assinatura como face de $\bTGL$); a
identificação canônica da matriz-S de fronteira
($\mathcal S_\partial=W_+^\dagger W_-$, espalhamento do cociclo de Connes)
e a hierarquia dos quatro operadores; a localização de $\sqrt e$ na entropia
de Araki; a articulação das duas condicionalidades (H)/(U) num só degrau
(Teorema~\ref{thm:degrau}); o Teorema da Escala (\S\ref{sec:escala}:
$\alpha(0)$ como ausência de escala --- congelamento IV, conteúdo da
fronteira, coisa julgada do fluxo wilsoniano, inscritibilidade
$\mathrm{III}_1$ --- objeção~d fechada condicional, com predição
falsificável própria); o Teorema da Terminalidade pela Verdade (T3 fechado:
Frigerio$+$Tomiyama$+$Takesaki$+$Gelfand, prova do Operador,
\S\ref{sec:terminalidade}); a evasão do no-go de Fewster--Verch pelo
colapso terminal; o Levantamento condicional na ordem de Jacobson
(\S\ref{sec:levantamento}); a Ordem Nominal (\S\ref{sec:ordemnominal}: a
covariância pela Lei comum --- Hadamard/microlocal, Radzikowski--Verch ---
não pelo juiz comum; as duas faces do Nome harmonizadas pela meia-nat);
a Geometricidade (\S\ref{sec:geometricidade}: ser $=$ ter geometria; a
geometria gera-se do dado modular, não se confere a ele --- a inversão do
ônus que dissolve o BW geral em coincidência [\textsf{REAL}] $+$
consistência de colagem [sombra]; a gramática do título como último
axioma da casa); o Teorema do Contínuo (montagem condicional,
\S\ref{sec:continuo}: Haagerup--BDF, classificação abeliana, Rieffel,
reconstrução de Connes, OS com positividade de reflexão \emph{automática}
--- a meia-nat fatora em duas quartas-medidas); os Três Fechos Técnicos
(\S\ref{sec:tresfechos}: a identidade do fluxo do radical
$e^{t\mathcal L}=\int V_s(\cdot)V_s^*d\nu_t$, $V_s=e^{is\sqrt K}$; T1
bem-posto no produto cruzado, o lar físico por CLPW; tripla-limite $=$
círculo canônico; levantamento GH $=$ truncamento espectral de
Connes--van Suijlekom); o Teorema do Rio (\S\ref{sec:rio}: Heráclito $=$
Lema T, outer sem retorno; Crátilo $=$ fluxo dos pesos trivial de
$\mathrm{III}_1$, o rio sem redemoinhos; pleno invariante vs.\ substância
outer; e $\mathrm{Fix}(\sigma_t)=\mathrm{Fix}(\Phi_t)=\mathcal C$ --- tudo
flui, exceto a verdade).

\smallskip\noindent\textbf{Confirmado na sombra (tipo I):} admissibilidade
algébrica de Bianchi da torção tracial; $\mathcal R_T$ rank-1 magnitude
$\bTGL$; $1/\bTGL=83$ fixo; Lindblad colapsa para atrator; núcleo 1-dim de
$\Pi_{\mathrm{bulk}}$; flecha do tempo via atrator único; $p=1\Leftrightarrow
S=0$; cociclo de Connes computado (unitariedade, regra de cadeia,
entrelaçamento a $\sim10^{-14}$); $\langle-i\,\dot u_0\rangle_\rho
=S_{\mathrm{Araki}}$ exato; código holográfico exatamente corretível
(Knill--Laflamme, resíduo $\sim10^{-15}$; fidelidade $1$ pela linha do ponto
único); covariância de $\mathcal R\circ\mathcal C$ e independência de
horizonte da fonte ($10^{-14}$--$10^{-17}$); lei de ressurreição
$1/(2\sqrt{\bTGL})$ (expoente $-0{,}51$); $\Phi_\infty=E$ a $10^{-27}$;
unicidade da fatoração (espaço nulo de dimensão zero); covariância
herdada $UEU^{-1}=E'$ a $10^{-16}$; em $M_n$ a valoração termina em $0=1$
(certificado), em $\mathcal C$ exatamente $n$ caracteres; rito-independência
do terminal dentro da ordem nominal ($10^{-27}$, cinco ritos); outra verdade
fora da ordem (controle negativo); transporte da classe ($10^{-12}$); pureza
monotônica não-crescente sob todo rito admissível; a geometria gerada é
métrica genuína (triangular a $10^{-18}$), local --- colagem exata, $0$
bit a bit --- e covariante ($10^{-16}$); o indefinido ($K=c\,\mathbb 1$)
gera taxas $\equiv0$, gerador nulo e $\Phi_t=\mathrm{id}$: nada não tem
contorno, Verbo, nem verdade; qualquer distinção gera geometria positiva;
a cascata de sombras é GH-Cauchy ao contínuo (taxa $0{,}500000$/geração,
$R^2=1{,}0000000$); o limite é canônico (independe do esquema de
refinamento) e carrega a medida; positividade de reflexão automática
(identidade $M=\mathrm{Gram}$ a $10^{-16}$, autovalores estritamente
positivos); a identidade do fluxo do radical a $4\times10^{-16}$; escada
de Dirac exata ($10^{-15}$), dimensão de Weyl $0{,}9944$, dualidade de
Connes ao contínuo ($\to1{,}9\times10^{-5}$); razão da testemunha ótima
truncada $0{,}9772\to0{,}9971$ ($O(1/M)$); a sombra perde os redemoinhos
($R$: $0{,}006\to0{,}029$ com $n$); pleno a $10^{-16}$ com substância
movendo $\geq83\%$; trajetória dissipativa injetiva; e
$\|P_{\mathrm{fluxo}}-P_{\mathrm{colapso}}\|=0$ exato.

\smallskip\noindent\textbf{Refutado (pelos números, registrado para não
retornar):} $\Lambda=\bTGL H_0^2$; $\thM=$ holonomia/Berry; $\mathcal R_T=$
torção (é simétrico); $T=N-\partial N$ (é zero); atrator único $\Rightarrow$
eixo temporal único; $g=g^{(+)}-g^{(-)}$ (subtração de tensores);
$\Pi_{\mathrm{bulk}}$ só-com-rotação gera tempo; $\sqrt e$ como autovalor
de $\mathcal S_\partial$ (unitário: espectro de fases) ou como razão
modular singularizada ($\mathrm{III}_1$: razão $=\mathbb R_+$); a Meia-Nat
como teorema algébrico (Araki sem \emph{gap}) ou como limiar universal de
detector (acompanha $\tau_{\det}$).

\smallskip\noindent\textbf{Aberto (programa, com prioridade):}
a escolha do Lindblad $\mathcal L=\sqrt\bTGL
\sqrt{K_\partial}$ dentre as $f(K_\partial)$ (\S\ref{sec:objecoes}c); a
reconstrução lorentziana de Connes sem traço, em regime $\mathrm{III}_1$
(Axioma~\ref{ax:lorentz}, hipótese H --- re-condicionada aos três axiomas
da casa, \S\ref{sec:axiomascasa}) --- campo matemático aberto, cujo
obstáculo central (ausência de traço/projetores em $\mathrm{III}_1$) a TGL
endereça via produto cruzado de Takesaki (\S\ref{sec:takesaki}); a derivação
(não a inscrição holográfica) da dimensão transversal $3$
(\S\ref{sec:tres}); a derivação explícita de $K_{\bTGL}$ da torção no
contínuo; a magnitude absoluta de $\Lambda$ (problema comum a toda a física); a
completude assintótica modular (existência dos limites de Møller $W_\pm$);
a transferência de $|\mathcal R|^2=\bTGL$ do modelo finito ao canal de
fronteira gravitacional; a Hipótese de Universalidade (U) da Face C --- o
levantamento covariante de $\mathcal R\circ\mathcal C$ a $\mathrm{III}_1$
genuína (a face dinâmica do mesmo degrau da reconstrução); a derivação da
escala $\tau_\star$ (dimensional: exige o corte UV/Planck). Após o fecho de
T3: o resíduo do Levantamento comprime-se a \textbf{(i)}~BW generalizado
não-perturbativo para horizontes não-Killing e \textbf{(ii)}~T1
(ergodicidade do colapso em $\mathrm{III}_1$ genuína) --- ambos nomeados,
com a ordem linear coberta condicionalmente (\S\ref{sec:levantamento}).
Após Ax.N (\S\ref{sec:ordemnominal}), (i) refina-se à forma mínima: a
\emph{realização geométrica} da classe outer --- já estado-independente e
geométrica pela ordem nominal --- além da ordem linear, em horizontes
não-Killing. Após Ax.G (\S\ref{sec:geometricidade}), o ônus inverte-se e
(i) dissolve-se em definição $+$ consistência (coincidência onde há
geometria pré-dada, \textsf{REAL}; colagem exata na sombra): resta
\textbf{T1} --- e a pergunta clássica do BW sobre geometrias de entrada,
registrada como problema externo que não bloqueia a cadeia. Forma final
dos condicionais do contínuo (\S\ref{sec:continuo}): \textbf{T1}
(ergodicidade em $\mathrm{III}_1$ genuína), os axiomas de Connes para a
tripla-limite (suportados pela Fractalização), e a extensão da
convergência GH a $\mathrm{III}_1$ (provada na sombra; conjecturada via
Rieffel) --- três itens de análise; nenhum de ontologia. Após os Três
Fechos (\S\ref{sec:tresfechos}): \textbf{nenhum item sem mecanismo} ---
resta a redação operador-algébrica completa (dois lemas padrão) e a
generalização dimensional por produtos: trabalho de escrita e montagem,
não de fundamento.

% =====================================================================
\section{Fechamos: a Síntese da TGL}
% =====================================================================
A Teoria da Gravitação Luminodinâmica fecha-se como \emph{uma cadeia única de
operadores}, governada por uma só constante singular não-derivável. Não é uma
coleção de analogias: é um gerador, lido em registros distintos.

\subsection{A cadeia, em uma sequência}
\begin{equation}
\underset{\text{axioma}}{g=\sqrt{|L_\varphi|}}
\;\to\;
\underset{\text{singularidade}}{\bTGL=\alpha e^{1/2}}
\;\to\;
\underset{\text{unificação}}{\mathcal O_\beta\mathcal L=\sqrt\bTGL\,\mathcal L}
\;\to\;
\underset{\text{gesto}}{S=J\Delta^{1/2}}
\end{equation}
\begin{equation}
\to\;
\underset{\text{jurisdição}}{T_t=e^{-t\mathcal L}}
\;\to\;
\underset{\text{tempo}}{\text{coisa julgada}}
\;\to\;
\underset{\text{espaço}}{\text{dobra da força}}
\;\to\;
\underset{\text{Ponte}}{(M^{1,3},g,\Gamma_{\mathrm{LC}}+K_{\bTGL})}
\end{equation}
\begin{equation}
\to\;
\underset{\text{entrelaçador}}{u_t=[D\rho:D\rho^\star]_t}
\;\to\;
\underset{\text{matriz-S}}{\mathcal S_\partial=W_+^\dagger W_-}
\;\to\;
\underset{\text{fonte}}{\mathcal P_{\mu\nu}[K_\partial]}
\;\to\;
\underset{\text{Face C}}{G_{\mu\nu}+\Lambda g_{\mu\nu}
=8\pi G\,\mathcal P_{\mu\nu}}.
\end{equation}
A primeira linha é o fecho cinemático (o palco); a segunda, o fecho dinâmico
(o drama). As duas condicionalidades --- (H) e (U) --- são um só degrau
(Teorema~\ref{thm:degrau}).

\subsection{O princípio único}
\begin{principle}[A Síntese]
Tudo na TGL é a inscrição de uma só distinção mínima contra a coincidência
total. O zero absoluto (a justiça, a morte, $1{=}0$, $p=1$) é inatingível e
inabitável. \textbf{Ser é não coincidir consigo} --- pagar $\bTGL$, a
meia-nat, o custo de inscrever a fronteira da distinção (a Palavra do
Juramento). Dessa única inscrição decorrem, como faces:
\begin{itemize}[noitemsep]
  \item a \emph{luz} (o autovetor que paga $\sqrt\bTGL$ por nível para viver);
  \item o \emph{tempo} (o irreversível, a coisa julgada, o eixo excluído);
  \item o \emph{espaço} (a dobra da força que não pode coincidir);
  \item a \emph{verdade} (pagar o custo) e a \emph{mentira} (fingir não pagar);
  \item a \emph{geometria identificável} (a jurisdição, contra a justiça
  não-identificável);
  \item a \emph{assinatura} $(1,3)$ (a face geométrica de $\bTGL$);
  \item a \emph{curvatura} (a resposta covariante conservada da geometria
  à reconstrução modular da fronteira: a Face C);
  \item a \emph{covariância} (a Ordem Nominal: a Lei comum sob a qual toda
  comarca julga --- o Nome que ordena por ser inalcançável);
  \item a \emph{geometricidade} (ser $=$ ter geometria: o espaço é causado
  pelo dado modular, não dado a ele --- e nada, sem contorno, não recebe
  artigo);
  \item o \emph{contínuo} (a sombra é densa nele --- hiperfinitude: cada
  validação finita é um ponto do limite; e a travessia ao lorentziano é
  positiva porque é um quadrado);
  \item o \emph{fluxo do radical} (a própria $g=\sqrt{|L|}$ executa o
  colapso: $V_s=e^{is\sqrt K}$, e o colapso é a média gaussiana das
  flexões da raiz);
  \item o \emph{rio} (nunca duas vezes, porque outer; nem uma vez, porque
  o espectro é contínuo; e tudo flui exceto a verdade ---
  $\mathrm{Fix}$ do tempo e do juízo coincidem em $\mathcal C$).
\end{itemize}
A unicidade de cada uma é a unicidade de $\bTGL$. A constante não se deriva
porque \emph{é} o que não se deriva: a singularidade postulada, o Nome, o
um-que-é-zero no registro do Nome (jamais no da álgebra).
\end{principle}

\subsection{A frase honesta}
A TGL reformula a gravidade, o tempo e o espaço como \emph{inscrição modular}
governada por uma constante singular não-derivável, $\bTGL=\alpha\sqrt e$. A
estrutura é fechada e coerente; a singularidade é demonstrada irredutível; a
reconstrução como variedade lorentziana suave é \emph{teorema condicional}
(hipótese H), e as equações de campo são \emph{teorema condicional}
(Hipótese de Universalidade U) --- duas faces do mesmo degrau: o
levantamento covariante global da estrutura modular ao contínuo
$\mathrm{III}_1$, programa em curso, com a pergunta precisa endereçada ao
campo. Não se afirma prova
incondicional: afirmar isso seria a mentira que a própria teoria define ---
pretender um custo não pago.

\begin{center}
\rule{0.65\textwidth}{0.4pt}\\[0.6em]
\small
$\bTGL=\alpha\sqrt e$ --- a singularidade que não se deriva por ser o que é.\\
Uma constante, um gerador, em todos os registros. Zero parâmetros livres.\\
O custo é negação; a mentira, a pretensão de não pagar; a verdade paga.\\
A justiça é o zero absoluto; a jurisdição (Lindblad) segura o identificável.\\
O tempo é a coisa julgada; o espaço, a dobra da força que não coincide.\\
A assinatura $(1,3)$ é a face de $\bTGL$: paga-se a meia-nat e exclui-se $1{=}0$.\\
A unicidade do tempo é a unicidade da singularidade --- ancorada, não derivada.\\
O custo do zero absoluto é ``haja luz''. Ser é pagar a meia-nat.\\[0.4em]
\textit{Ao meu Deus: não o justo, mas o misericordioso, e por isso o
justificador.}\\
\textit{A completude é o contorno do que é bastante.}\\[0.3em]
\textbf{TETELESTAI} --- consumado como reconhecido, não como provado a partir de.
\end{center}

\end{document}